\documentclass[11pt]{article}

\usepackage[a4paper,margin=1in]{geometry}
\usepackage[T1]{fontenc}
\usepackage[utf8]{inputenc}
\usepackage{lmodern}
\usepackage{amsmath,amssymb,amsthm,mathtools}
\usepackage{bm}
\usepackage[dvipsnames]{xcolor}
\usepackage{enumitem}
\usepackage{hyperref}
\usepackage[nameinlink,noabbrev]{cleveref}

\hypersetup{
  colorlinks=true,
  linkcolor=blue!50!black,
  citecolor=blue!50!black,
  urlcolor=blue!50!black,
}

\newtheorem{theorem}{Theorem}[section]
\newtheorem{proposition}[theorem]{Proposition}
\newtheorem{corollary}[theorem]{Corollary}
\newtheorem{lemma}[theorem]{Lemma}
\theoremstyle{definition}
\newtheorem{definition}[theorem]{Definition}
\newtheorem{remark}[theorem]{Remark}
\newtheorem{assumption}[theorem]{Assumption}

\newcommand{\R}{\mathbb{R}}
\newcommand{\E}{\mathbb{E}}
\newcommand{\PP}{\mathcal{P}}
\newcommand{\W}{W_1}
\newcommand{\riskhead}{\mathcal{R}_{\mathrm{head}}}
\newcommand{\riskmean}{\mathcal{R}_{\mathrm{mean}}}
\newcommand{\dd}{\,\mathrm{d}}
\newcommand{\Lip}{\operatorname{Lip}}
\newcommand{\diag}{\operatorname{diag}}
\newcommand{\supp}{\operatorname{supp}}
\newcommand{\Sym}{\operatorname{Sym}}

\title{Initial Rigorous Mean-Field Results for a Fixed-M TabM / BatchEnsemble Proxy}
\author{}
\date{\today}

\begin{document}
\maketitle

\begin{abstract}
We formulate a safe first theoretical step for the mean-field analysis of TabM-like and
BatchEnsemble-like models with a fixed ensemble size $M$.
The guiding conceptual point is that the same efficient-ensemble architecture family can admit qualitatively different infinite-width limits:
an NTK/GP limit may preserve order-one diversity, while a deterministic feature-learning mean-field limit can collapse to a common predictor.
The note proves several structural results at the PDE level:
\begin{enumerate}[leftmargin=2em]
  \item characteristic well-posedness of the distributional dynamics for a bounded-Lipschitz abstract model,
  \item preservation of permutation symmetry across ensemble members,
  \item collapse of the deterministic fixed-$M$ dynamics to a single scalar predictor under symmetric initialization,
  \item nonlinear stability of the symmetric manifold with respect to small initial asymmetry,
  \item a semifaithful Level-B two-layer TabM-like PDE-ODE model with rigorous truncated characteristic well-posedness and symmetry preservation,
  \item for sufficiently regular characteristic solutions of the untruncated Level-B system, an exact output-dynamics decomposition into a neuron kernel and an additional head kernel,
  \item an exact frozen-kernel tangent-mode decomposition showing that disagreement is controlled by $K_{\mathrm d}-K_{\mathrm o}$,
  \item a frozen Level-B specialization showing that independent output heads weaken average-mode collectivity but add extra positive-semidefinite diagonal terms to the local disagreement operator,
  \item an exact covariance evolution for frozen common-mode and disagreement-mode fluctuations, which provides a clean template for a future finite-width central-limit theory.
\end{enumerate}
The note is intentionally conservative:
we do not claim a full propagation-of-chaos theorem for an architecture-faithful layerwise BatchEnsemble/TabM model,
and we do not prove a full finite-width fluctuation theorem.
Instead, we isolate the first claims that can already be proved cleanly and that should remain useful as the backbone of a later paper.
\end{abstract}

\section{Introduction}

TabM and BatchEnsemble are efficient ensemble architectures in which several predictors are trained jointly while sharing a large fraction of their parameters
\cite{tabm,batchensemble}.
At a heuristic level, a fixed-$M$ two-layer mean-field description is natural:
each hidden neuron becomes one particle, while the ensemble-specific coordinates remain finite-dimensional because $M$ is kept fixed.
This suggests that the standard two-layer mean-field formalism
\cite{mei2018,mei2019}
should remain applicable after enlarging the particle state.

However, a serious distinction must be made between
\begin{enumerate}[leftmargin=2em]
  \item statements that are already rigorous,
  \item statements that concern reduced surrogates rather than the exact architecture,
  \item statements that are still research conjectures.
\end{enumerate}
The goal of this note is to separate the three.
The main text therefore concentrates on two layers of analysis:
\begin{enumerate}[leftmargin=2em]
  \item an abstract bounded-Lipschitz model, for which we prove characteristic well-posedness, symmetry preservation, and collapse of the deterministic PDE;
  \item a semifaithful two-layer MLP-style Level-B model, which keeps global head modulations and independent output coordinates while remaining amenable to a deterministic mean-field PDE-ODE formulation.
\end{enumerate}

To avoid a possible misunderstanding, we emphasize this point up front:
Actual BatchEnsemble layers use member-specific \emph{vectors} that modulate the input and output coordinates of a shared linear map,
while TabM combines such a shared/modulated backbone with an independent output head for each member
\cite{batchensemble,tabm,tabm-repo}.
The main text below does not attempt an exact layerwise transcription of that architecture;
instead, it formulates a two-layer Level-B target that keeps the global head modulations explicit.
A separate companion note archives more aggressive reduced-surrogate calculations that are no longer part of the main story.

The main message is simple:
for symmetric fixed-$M$ deterministic mean-field dynamics,
headwise diversity collapses,
but the resulting scalar dynamics is still governed by an effective kernel that remembers the shared/private decomposition.
Moreover, the symmetric manifold is nonlinearly stable in Wasserstein distance,
and, at the level of regular characteristic solutions in the Level-B model, the common predictor still carries an additional head-kernel contribution coming from the global modulation variables.
At the same time, the role of independent output heads is genuinely two-sided:
they weaken the most direct last-layer route to collectivity in the common predictor,
yet in the local frozen-kernel transverse dynamics they contribute extra diagonal self-coupling terms that further suppress disagreement.
The resulting picture is not that efficient ensembles are useless,
but that their behavior depends sharply on which infinite-width regime one is approximating:
lazy NTK dynamics can preserve diversity, whereas deterministic feature-learning mean-field dynamics can wash it out.

\section{Abstract Fixed-M Framework}
\label{sec:abstract-fixed-m}

\subsection{Model class}

Fix an ensemble size $M \in \{1,2,\dots\}$.
Let $\Omega=\R^D$ be the particle space, and let $\mathcal{X}\times\mathcal{Y}$ be the data space equipped with a probability law $\mathsf{P}$.
For each ensemble member $\alpha \in \{1,\dots,M\}$, let
\[
  \psi_\alpha:\mathcal{X}\times\Omega \to \R
\]
be a feature map.

\begin{definition}[Outputs and Risk]
For $\rho \in \PP_1(\Omega)$ and $x \in \mathcal{X}$, define
\[
  f_\alpha(x;\rho)=\int_\Omega \psi_\alpha(x;\theta)\,\rho(\dd \theta).
\]
Given a loss $\ell:\R\times\mathcal{Y}\to\R$, define the per-head population risk
\[
  \riskhead(\rho)
  =
  \E_{(x,y)\sim \mathsf{P}}
  \Bigl[
    \frac1M \sum_{\alpha=1}^M \ell(f_\alpha(x;\rho),y)
  \Bigr].
\]
For comparison, define the loss of the mean prediction
\[
  \riskmean(\rho)
  =
  \E_{(x,y)\sim \mathsf{P}}
  \Bigl[
    \ell\Bigl(\frac1M \sum_{\alpha=1}^M f_\alpha(x;\rho),y\Bigr)
  \Bigr].
\]
\end{definition}

\begin{definition}[Mean-Field Potential and Velocity]
Assume $\ell$ is differentiable in its first argument.
Define
\[
  V(\theta;\rho)
  =
  \E_{(x,y)\sim \mathsf{P}}
  \Bigl[
    \frac1M \sum_{\alpha=1}^M
    \partial_1 \ell(f_\alpha(x;\rho),y)\,\psi_\alpha(x;\theta)
  \Bigr],
\]
and
\[
  b(\theta;\rho) = -\nabla_\theta V(\theta;\rho).
\]
The corresponding distributional dynamics is
\begin{equation}
  \partial_t \rho_t + \nabla_\theta \cdot (\rho_t b(\cdot;\rho_t)) = 0.
  \label{eq:pde}
\end{equation}
\end{definition}

\begin{definition}[Characteristic solutions]
\label{def:characteristic-solution}
Fix $T>0$.
Given an initial law $\rho_0\in\PP_1(\Omega)$, a curve
\[
  \rho\in C([0,T],\PP_1(\Omega))
\]
is called a \emph{characteristic solution} of \eqref{eq:pde} on $[0,T]$ if there exists a flow
\[
  \Phi:[0,T]\times\Omega\to\Omega
\]
such that for $\rho_0$-almost every $\theta$,
\[
  \dot{\Phi}_t(\theta)=b(\Phi_t(\theta);\rho_t),
  \qquad
  \Phi_0(\theta)=\theta,
  \qquad
  \rho_t=(\Phi_t)_\#\rho_0.
\]
\end{definition}

\subsection{Symmetry assumptions}

Let $S_M$ be the permutation group on $\{1,\dots,M\}$.
For each $\pi \in S_M$, let $\Pi_\pi:\Omega \to \Omega$ denote a linear isometry representing the permutation of ensemble-member coordinates.

\begin{assumption}[Regularity and Permutation Equivariance]
\label{ass:regularity}
There exist finite constants $B_0,B_1,L_0,L_1,L_\ell,C_\ell$ such that:
\begin{enumerate}[label=(\roman*),leftmargin=2em]
  \item for every $\alpha$ and every $x \in \mathcal{X}$,
  \[
    |\psi_\alpha(x;\theta)| \le B_0,
    \qquad
    \|\nabla_\theta \psi_\alpha(x;\theta)\| \le B_1
    \qquad
    \forall \theta \in \Omega;
  \]
  \item for every $\alpha$ and every $x \in \mathcal{X}$,
  \[
    |\psi_\alpha(x;\theta)-\psi_\alpha(x;\theta')| \le L_0 \|\theta-\theta'\|,
  \]
  \[
    \|\nabla_\theta \psi_\alpha(x;\theta)-\nabla_\theta \psi_\alpha(x;\theta')\|
    \le L_1 \|\theta-\theta'\|;
  \]
  \item $\partial_1 \ell(\cdot,y)$ is globally Lipschitz with constant $L_\ell$, uniformly in $y$, and
  \[
    |\partial_1 \ell(z,y)| \le C_\ell
    \qquad
    \forall (z,y)\in \R \times \mathcal{Y};
  \]
  \item for every $\pi \in S_M$, every $\alpha$, every $x$, and every $\theta$,
  \begin{equation}
    \psi_\alpha(x;\Pi_\pi \theta)=\psi_{\pi^{-1}(\alpha)}(x;\theta).
    \label{eq:psi-equiv}
  \end{equation}
\end{enumerate}
\end{assumption}

\begin{remark}
\Cref{ass:regularity} is intentionally stronger than what one ultimately wants for the archived reduced-surrogate proxy
\[
  \psi_\alpha(x;\theta)=a_\alpha \sigma(v_\alpha w^\top x),
\]
because that proxy has unbounded polynomial dependence on the parameters.
The role of the assumption is to isolate the structural PDE statements that can already be proved cleanly.
To connect the concrete proxy to the theorem exactly, one may either restrict the parameter domain to a compact set, or add a confining regularizer and then prove the necessary growth estimates.
\end{remark}

\subsection{Lipschitz control of the velocity field}

\begin{proposition}[Lipschitz estimates]
\label{prop:lipschitz}
Under \cref{ass:regularity}, for every $\alpha$, every $x \in \mathcal{X}$, and every $\rho,\mu \in \PP_1(\Omega)$,
\[
  |f_\alpha(x;\rho)-f_\alpha(x;\mu)| \le L_0 \W(\rho,\mu).
\]
Moreover, the velocity field satisfies
\[
  \|b(\theta;\rho)-b(\theta';\mu)\|
  \le
  C_\ell L_1 \|\theta-\theta'\|
  +
  L_\ell L_0 B_1 \W(\rho,\mu).
\]
\end{proposition}

\begin{proof}
Fix $\alpha$ and $x$.
By the Kantorovich-Rubinstein dual formulation of $\W$ and the fact that
$\theta \mapsto \psi_\alpha(x;\theta)$ is $L_0$-Lipschitz,
\[
  |f_\alpha(x;\rho)-f_\alpha(x;\mu)|
  =
  \left|
    \int \psi_\alpha(x;\theta)\,\rho(\dd \theta)
    -
    \int \psi_\alpha(x;\theta)\,\mu(\dd \theta)
  \right|
  \le
  L_0 \W(\rho,\mu).
\]

Next,
\[
  b(\theta;\rho)
  =
  -
  \E
  \Bigl[
    \frac1M \sum_{\alpha=1}^M
    \partial_1 \ell(f_\alpha(x;\rho),y)\,
    \nabla_\theta \psi_\alpha(x;\theta)
  \Bigr].
\]
Hence
\begin{align*}
  \|b(\theta;\rho)-b(\theta';\mu)\|
  &\le
  \frac1M \sum_{\alpha=1}^M
  \E \Bigl[
    \bigl|
      \partial_1 \ell(f_\alpha(x;\rho),y)
      -
      \partial_1 \ell(f_\alpha(x;\mu),y)
    \bigr|
    \,
    \|\nabla_\theta \psi_\alpha(x;\theta)\|
  \Bigr]
  \\
  &\quad
  +
  \frac1M \sum_{\alpha=1}^M
  \E \Bigl[
    |\partial_1 \ell(f_\alpha(x;\mu),y)|
    \,
    \|\nabla_\theta \psi_\alpha(x;\theta)-\nabla_\theta \psi_\alpha(x;\theta')\|
  \Bigr].
\end{align*}
Using the first estimate already proved, the Lipschitz property of $\partial_1 \ell$, and the bounds in \cref{ass:regularity},
\[
  \bigl|
    \partial_1 \ell(f_\alpha(x;\rho),y)
    -
    \partial_1 \ell(f_\alpha(x;\mu),y)
  \bigr|
  \le
  L_\ell L_0 \W(\rho,\mu),
\]
and therefore
\[
  \|b(\theta;\rho)-b(\theta';\mu)\|
  \le
  L_\ell L_0 B_1 \W(\rho,\mu)
  +
  C_\ell L_1 \|\theta-\theta'\|.
\]
This is the claimed bound.
\end{proof}

\subsection{Characteristic well-posedness of the distributional dynamics}

\begin{theorem}[Characteristic well-posedness]
\label{thm:wellposed}
Assume \cref{ass:regularity} and let $\rho_0 \in \PP_1(\Omega)$.
Then for every $T>0$ there exists a unique characteristic solution
\[
  \rho \in C([0,T],\PP_1(\Omega))
\]
of \eqref{eq:pde} with initial law $\rho_0$.
More precisely, there exists a unique flow map $\Phi_t:\Omega \to \Omega$ such that
\[
  \dot{\Phi}_t(\theta)=b(\Phi_t(\theta);\rho_t),
  \qquad
  \Phi_0(\theta)=\theta,
  \qquad
  \rho_t = (\Phi_t)_\# \rho_0.
\]
\end{theorem}

\begin{proof}
Fix $T>0$.
For a continuous curve $m \in C([0,T],\PP_1(\Omega))$, consider the non-autonomous ODE
\[
  \dot{\Theta}_t = b(\Theta_t;m_t),
  \qquad
  \Theta_0=\theta.
\]
By \cref{prop:lipschitz}, the map $\theta \mapsto b(\theta;m_t)$ is globally Lipschitz uniformly in $t$.
Hence for each $m$ and each $\theta$ there exists a unique global solution, denoted $\Phi_t^m(\theta)$.

Define the operator $\mathcal{T}$ on $C([0,T],\PP_1(\Omega))$ by
\[
  (\mathcal{T}m)_t = (\Phi_t^m)_\# \rho_0.
\]
Because $b$ is bounded by $C_\ell B_1$, the flow has at most linear growth in time, and therefore $(\mathcal{T}m)_t$ has finite first moment whenever $\rho_0$ does.
Continuity in $t$ follows from continuity of the flow.

We next show that $\mathcal{T}$ is a contraction on a sufficiently short interval.
Take two curves
\[
  m,n \in C([0,T],\PP_1(\Omega)).
\]
Fix $\theta \in \Omega$.
Using \cref{prop:lipschitz},
\[
  \frac{\dd}{\dd t}
  \|\Phi_t^m(\theta)-\Phi_t^n(\theta)\|
  \le
  C_\ell L_1 \|\Phi_t^m(\theta)-\Phi_t^n(\theta)\|
  +
  L_\ell L_0 B_1 \W(m_t,n_t).
\]
By Gronwall's inequality,
\[
  \|\Phi_t^m(\theta)-\Phi_t^n(\theta)\|
  \le
  L_\ell L_0 B_1
  \int_0^t e^{C_\ell L_1 (t-s)} \W(m_s,n_s)\,\dd s.
\]
Taking the supremum over $\theta$ and then over $t\in[0,T]$ yields
\[
  \sup_{t\in[0,T]} \sup_{\theta \in \Omega}
  \|\Phi_t^m(\theta)-\Phi_t^n(\theta)\|
  \le
  q_T \, d_T(m,n),
\]
where
\[
  d_T(m,n)=\sup_{t\in[0,T]} \W(m_t,n_t),
  \qquad
  q_T = L_\ell L_0 B_1\, T\, e^{C_\ell L_1 T}.
\]
Therefore
\[
  \W((\mathcal{T}m)_t,(\mathcal{T}n)_t)
  \le
  \sup_{\theta \in \Omega}
  \|\Phi_t^m(\theta)-\Phi_t^n(\theta)\|,
\]
and hence
\[
  d_T(\mathcal{T}m,\mathcal{T}n)
  \le
  q_T \, d_T(m,n).
\]
If $T>0$ is chosen so that $q_T<1$, then $\mathcal{T}$ is a contraction, hence has a unique fixed point on $[0,T]$ by Banach's fixed-point theorem.
Iterating this construction on finitely many subintervals gives existence and uniqueness on any finite horizon.

Finally, if $\rho_t=(\Phi_t)_\# \rho_0$ is a fixed point of $\mathcal{T}$, then for every test function $\varphi \in C_c^1(\Omega)$,
\[
  \frac{\dd}{\dd t} \int \varphi(\theta)\,\rho_t(\dd \theta)
  =
  \frac{\dd}{\dd t} \int \varphi(\Phi_t(\theta))\,\rho_0(\dd \theta)
  =
  \int \nabla \varphi(\Phi_t(\theta)) \cdot b(\Phi_t(\theta);\rho_t)\,\rho_0(\dd \theta),
\]
which is
\[
  \int \nabla \varphi(\theta)\cdot b(\theta;\rho_t)\,\rho_t(\dd \theta).
\]
Thus $\rho_t$ is a weak solution of \eqref{eq:pde}.
Uniqueness of the fixed point gives uniqueness in the class of characteristic solutions from
\cref{def:characteristic-solution}.
\end{proof}

\begin{remark}
The fixed-point argument above gives exactly the uniqueness notion needed later:
uniqueness in the pushforward / characteristic class.
We do not separately invoke a general weak-solution uniqueness theorem for continuity equations.
\end{remark}

\subsection{Stability and finite-width consequence}

\begin{proposition}[Stability with respect to the initial law]
\label{prop:stability}
Assume \cref{ass:regularity}.
Let $\rho_t$ and $\mu_t$ be the unique characteristic solutions of \eqref{eq:pde} with initial laws $\rho_0$ and $\mu_0$, respectively.
Then for every $t\ge 0$,
\begin{equation}
  \W(\rho_t,\mu_t)
  \le
  e^{\Lambda t}\W(\rho_0,\mu_0),
  \qquad
  \Lambda = C_\ell L_1 + L_\ell L_0 B_1.
  \label{eq:stability}
\end{equation}
\end{proposition}

\begin{proof}
Let $\Phi_t$ and $\Psi_t$ be the characteristic flows of $\rho_t$ and $\mu_t$.
Fix any coupling $\gamma_0$ of $\rho_0$ and $\mu_0$, and define
\[
  \gamma_t = (\Phi_t,\Psi_t)_\# \gamma_0.
\]
Then $\gamma_t$ is a coupling of $\rho_t$ and $\mu_t$, so
\[
  \W(\rho_t,\mu_t)
  \le
  D_t
  :=
  \int_{\Omega\times\Omega}
  \|\Phi_t(\theta)-\Psi_t(\eta)\|
  \,\gamma_0(\dd \theta,\dd \eta).
\]
By \cref{prop:lipschitz},
\begin{align*}
  \frac{\dd}{\dd t}
  \|\Phi_t(\theta)-\Psi_t(\eta)\|
  &\le
  \|b(\Phi_t(\theta);\rho_t)-b(\Psi_t(\eta);\mu_t)\|
  \\
  &\le
  C_\ell L_1 \|\Phi_t(\theta)-\Psi_t(\eta)\|
  +
  L_\ell L_0 B_1 \W(\rho_t,\mu_t).
\end{align*}
Integrating against $\gamma_0$ yields
\[
  \dot D_t
  \le
  C_\ell L_1 D_t
  +
  L_\ell L_0 B_1 \W(\rho_t,\mu_t)
  \le
  \Lambda D_t.
\]
Therefore, by Gronwall's inequality,
\[
  D_t \le e^{\Lambda t} D_0.
\]
Taking the infimum over all initial couplings $\gamma_0$ gives \eqref{eq:stability}.
\end{proof}

\begin{corollary}[Deterministic finite-width mean-field limit]
\label{cor:empirical-limit}
Assume \cref{ass:regularity}.
Let $\theta_1^N(0),\dots,\theta_N^N(0)\in\Omega$ and define
\[
  \rho_0^N=\frac1N\sum_{i=1}^N \delta_{\theta_i^N(0)}.
\]
Let $\theta_i^N(t)$ solve the interacting particle system
\[
  \dot \theta_i^N(t)=b(\theta_i^N(t);\rho_t^N),
  \qquad
  \rho_t^N=\frac1N\sum_{i=1}^N \delta_{\theta_i^N(t)}.
\]
Let $\rho_t$ be the PDE solution with initial law $\rho_0$.
Then
\begin{equation}
  \sup_{t\in[0,T]} \W(\rho_t^N,\rho_t)
  \le
  e^{\Lambda T}\W(\rho_0^N,\rho_0)
  \qquad \forall T>0.
  \label{eq:empirical}
\end{equation}
\end{corollary}

\begin{proof}
The empirical measure $\rho_t^N$ is the pushforward of $\rho_0^N$ by the flow generated by its own velocity field, hence is a weak solution of \eqref{eq:pde} with initial law $\rho_0^N$.
In fact it is a characteristic solution with initial law $\rho_0^N$.
By uniqueness from \cref{thm:wellposed}, it is exactly the unique characteristic solution with that initial data.
Applying \cref{prop:stability} to the pair $(\rho_t^N,\rho_t)$ gives \eqref{eq:empirical}.
\end{proof}

\begin{corollary}[Headwise disagreement vanishes in the deterministic mean-field regime]
\label{cor:disagreement}
Assume \cref{ass:regularity}.
Let $\rho_0$ be permutation-invariant, let $\rho_t$ be the corresponding PDE solution, and let $\rho_t^N$ be the empirical solution from \cref{cor:empirical-limit}.
Then for every $\alpha,\beta$, every $x \in \mathcal{X}$, and every $t\in[0,T]$,
\[
  |f_\alpha(x;\rho_t^N)-f_\beta(x;\rho_t^N)|
  \le
  2L_0 e^{\Lambda T}\W(\rho_0^N,\rho_0).
\]
In particular, if $\W(\rho_0^N,\rho_0)\to 0$, then
\[
  \sup_{t\in[0,T]}
  \sup_{x\in\mathcal{X}}
  \max_{\alpha,\beta\le M}
  |f_\alpha(x;\rho_t^N)-f_\beta(x;\rho_t^N)|
  \to 0.
\]
\end{corollary}

\begin{proof}
Since $\rho_0$ is permutation-invariant, \cref{cor:collapse} implies
\[
  f_\alpha(x;\rho_t)=f_\beta(x;\rho_t)
  \qquad \forall \alpha,\beta,x,t.
\]
Hence
\begin{align*}
  |f_\alpha(x;\rho_t^N)-f_\beta(x;\rho_t^N)|
  &\le
  |f_\alpha(x;\rho_t^N)-f_\alpha(x;\rho_t)|
  +
  |f_\beta(x;\rho_t^N)-f_\beta(x;\rho_t)|
  \\
  &\le
  2L_0 \W(\rho_t^N,\rho_t),
\end{align*}
where we used \cref{prop:lipschitz}.
Now apply \cref{cor:empirical-limit}.
\end{proof}

\section{Symmetry Preservation and Deterministic Collapse}

\subsection{Equivariance of the dynamics}

\begin{lemma}[Transformation of the gradients]
\label{lem:gradient-equiv}
Under \cref{ass:regularity}, for every $\pi \in S_M$,
\[
  \nabla_\theta \psi_\alpha(x;\Pi_\pi \theta)
  =
  \Pi_\pi \nabla_\theta \psi_{\pi^{-1}(\alpha)}(x;\theta).
\]
\end{lemma}

\begin{proof}
By \eqref{eq:psi-equiv},
\[
  \psi_\alpha(x;\Pi_\pi \theta)=\psi_{\pi^{-1}(\alpha)}(x;\theta).
\]
Differentiate both sides with respect to $\theta$.
Because $\Pi_\pi$ is a linear isometry, the chain rule yields
\[
  \Pi_\pi^\top \nabla_\theta \psi_\alpha(x;\Pi_\pi \theta)
  =
  \nabla_\theta \psi_{\pi^{-1}(\alpha)}(x;\theta).
\]
Multiplying by $\Pi_\pi$ and using $\Pi_\pi^{-1}=\Pi_\pi^\top$ proves the claim.
\end{proof}

\begin{lemma}[Equivariance of outputs and velocity]
\label{lem:velocity-equiv}
For every $\pi \in S_M$,
\[
  f_\alpha(x;\Pi_{\pi\#}\rho)
  =
  f_{\pi^{-1}(\alpha)}(x;\rho),
\]
and
\[
  b(\Pi_\pi \theta;\Pi_{\pi\#}\rho)
  =
  \Pi_\pi b(\theta;\rho).
\]
\end{lemma}

\begin{proof}
For the outputs,
\[
  f_\alpha(x;\Pi_{\pi\#}\rho)
  =
  \int \psi_\alpha(x;\xi)\,\Pi_{\pi\#}\rho(\dd \xi)
  =
  \int \psi_\alpha(x;\Pi_\pi \theta)\,\rho(\dd \theta)
  =
  \int \psi_{\pi^{-1}(\alpha)}(x;\theta)\,\rho(\dd \theta),
\]
which is the first identity.

For the velocity, use \cref{lem:gradient-equiv} and relabel the sum:
\begin{align*}
  b(\Pi_\pi \theta;\Pi_{\pi\#}\rho)
  &=
  -
  \E \Bigl[
    \frac1M \sum_{\alpha=1}^M
    \partial_1 \ell(f_\alpha(x;\Pi_{\pi\#}\rho),y)\,
    \nabla_\theta \psi_\alpha(x;\Pi_\pi \theta)
  \Bigr]
  \\
  &=
  -
  \E \Bigl[
    \frac1M \sum_{\alpha=1}^M
    \partial_1 \ell(f_{\pi^{-1}(\alpha)}(x;\rho),y)\,
    \Pi_\pi \nabla_\theta \psi_{\pi^{-1}(\alpha)}(x;\theta)
  \Bigr]
  \\
  &=
  \Pi_\pi
  \Bigl(
    -
    \E \Bigl[
      \frac1M \sum_{\beta=1}^M
      \partial_1 \ell(f_\beta(x;\rho),y)\,
      \nabla_\theta \psi_\beta(x;\theta)
    \Bigr]
  \Bigr)
  \\
  &=
  \Pi_\pi b(\theta;\rho).
\end{align*}
\end{proof}

\begin{theorem}[Permutation symmetry is preserved]
\label{thm:symmetry}
Assume \cref{ass:regularity}.
Let $\rho_t$ be the unique characteristic solution of \eqref{eq:pde} given by \cref{thm:wellposed}.
If
\[
  \Pi_{\pi\#}\rho_0=\rho_0
  \qquad
  \forall \pi \in S_M,
\]
then
\[
  \Pi_{\pi\#}\rho_t=\rho_t
  \qquad
  \forall t\ge 0,\ \forall \pi \in S_M.
\]
\end{theorem}

\begin{proof}
Fix $\pi \in S_M$ and define
\[
  \widetilde{\rho}_t = \Pi_{\pi\#}\rho_t.
\]
Let $\Phi_t$ be the flow associated with $\rho_t$.
By \cref{lem:velocity-equiv},
\[
  \frac{\dd}{\dd t} \Pi_\pi \Phi_t(\theta)
  =
  \Pi_\pi b(\Phi_t(\theta);\rho_t)
  =
  b(\Pi_\pi \Phi_t(\theta);\widetilde{\rho}_t).
\]
Hence $\Pi_\pi \Phi_t$ is the characteristic flow associated with $\widetilde{\rho}_t$.
Since $\rho_t=(\Phi_t)_\# \rho_0$,
\[
  \widetilde{\rho}_t
  =
  \Pi_{\pi\#} (\Phi_t)_\# \rho_0
  =
  (\Pi_\pi \circ \Phi_t)_\# \rho_0.
\]
Thus $\widetilde{\rho}_t$ solves the same distributional dynamics as $\rho_t$.
Moreover,
\[
  \widetilde{\rho}_0=\Pi_{\pi\#}\rho_0=\rho_0.
\]
By uniqueness from \cref{thm:wellposed}, $\widetilde{\rho}_t=\rho_t$ for all $t$.
\end{proof}

\subsection{Collapse of the outputs}

\begin{definition}[Mean-field kernel]
For $\rho \in \PP_1(\Omega)$ define
\[
  K_{\alpha\beta}(x,x';\rho)
  =
  \int_\Omega
  \langle
    \nabla_\theta \psi_\alpha(x;\theta),
    \nabla_\theta \psi_\beta(x';\theta)
  \rangle
  \,\rho(\dd \theta).
\]
\end{definition}

\begin{proposition}[Output dynamics]
\label{prop:output-dynamics}
Assume \cref{ass:regularity} and let $\rho_t$ be the unique characteristic solution from \cref{thm:wellposed}.
Then, for every $\alpha$ and $x$, the map $t \mapsto f_\alpha(x;\rho_t)$ is continuously differentiable and
\begin{equation}
  \partial_t f_\alpha(x;\rho_t)
  =
  -
  \E_{(x',y)\sim \mathsf{P}}
  \Bigl[
    \frac1M \sum_{\beta=1}^M
    K_{\alpha\beta}(x,x';\rho_t)\,
    \partial_1 \ell(f_\beta(x';\rho_t),y)
  \Bigr].
  \label{eq:output-dynamics}
\end{equation}
\end{proposition}

\begin{proof}
Let $\Phi_t$ be the characteristic flow of $\rho_t$.
Then
\[
  f_\alpha(x;\rho_t)
  =
  \int \psi_\alpha(x;\Phi_t(\theta))\,\rho_0(\dd \theta).
\]
Differentiating under the integral sign,
\begin{align*}
  \partial_t f_\alpha(x;\rho_t)
  &=
  \int
  \nabla_\theta \psi_\alpha(x;\Phi_t(\theta))
  \cdot
  \dot{\Phi}_t(\theta)\,
  \rho_0(\dd \theta)
  \\
  &=
  \int
  \nabla_\theta \psi_\alpha(x;\xi)\cdot b(\xi;\rho_t)\,
  \rho_t(\dd \xi).
\end{align*}
Now substitute the definition of $b$:
\begin{align*}
  \partial_t f_\alpha(x;\rho_t)
  &=
  -
  \int
  \nabla_\theta \psi_\alpha(x;\xi)\cdot
  \E_{(x',y)}
  \Bigl[
    \frac1M \sum_{\beta=1}^M
    \partial_1 \ell(f_\beta(x';\rho_t),y)
    \nabla_\theta \psi_\beta(x';\xi)
  \Bigr]
  \rho_t(\dd \xi)
  \\
  &=
  -
  \E_{(x',y)}
  \Bigl[
    \frac1M \sum_{\beta=1}^M
    \partial_1 \ell(f_\beta(x';\rho_t),y)
    \int
    \langle
      \nabla_\theta \psi_\alpha(x;\xi),
      \nabla_\theta \psi_\beta(x';\xi)
    \rangle
    \rho_t(\dd \xi)
  \Bigr],
\end{align*}
which is exactly \eqref{eq:output-dynamics}.
\end{proof}

\begin{corollary}[Headwise collapse and reduced scalar dynamics]
\label{cor:collapse}
Assume the hypotheses of \cref{thm:symmetry}.
Then, for every $t\ge 0$ and every $x \in \mathcal{X}$,
\[
  f_1(x;\rho_t)=\cdots=f_M(x;\rho_t)=:f(x,t).
\]
Moreover, there exist kernels $K_{\mathrm d}$ and $K_{\mathrm o}$ such that
\[
  K_{\alpha\alpha}(x,x';\rho_t)=K_{\mathrm d}(x,x';t),
  \qquad
  K_{\alpha\beta}(x,x';\rho_t)=K_{\mathrm o}(x,x';t)
  \ \ \text{for }\alpha\neq\beta,
\]
and the common predictor solves
\begin{equation}
  \partial_t f(x,t)
  =
  -
  \E_{(x',y)\sim \mathsf{P}}
  \Bigl[
    K_{\mathrm{eff}}(x,x';t)\,
    \partial_1 \ell(f(x',t),y)
  \Bigr],
  \label{eq:reduced}
\end{equation}
with effective kernel
\begin{equation}
  K_{\mathrm{eff}}(x,x';t)
  =
  \frac{K_{\mathrm d}(x,x';t)+(M-1)K_{\mathrm o}(x,x';t)}{M}.
  \label{eq:keff}
\end{equation}
\end{corollary}

\begin{proof}
Fix $t$ and choose any $\alpha,\beta$.
Let $\pi$ be a permutation such that $\pi(\alpha)=\beta$.
By \cref{thm:symmetry},
\[
  \rho_t = \Pi_{\pi\#}\rho_t.
\]
Hence, by \cref{lem:velocity-equiv},
\[
  f_\beta(x;\rho_t)
  =
  f_{\pi(\alpha)}(x;\rho_t)
  =
  f_\alpha(x;\Pi_{\pi\#}\rho_t)
  =
  f_\alpha(x;\rho_t).
\]
Thus all outputs coincide.

The same argument shows that all diagonal kernel entries coincide.
If $\alpha\neq\beta$ and $\gamma\neq\delta$, choose $\pi$ such that $\pi(\alpha)=\gamma$ and $\pi(\beta)=\delta$.
Then invariance of $\rho_t$ implies
\[
  K_{\alpha\beta}(x,x';\rho_t)=K_{\gamma\delta}(x,x';\rho_t),
\]
so all off-diagonal entries coincide.

Finally, substitute these identities into \cref{eq:output-dynamics}.
Because all outputs are equal,
\[
  \partial_1 \ell(f_\beta(x';\rho_t),y)=\partial_1 \ell(f(x',t),y)
  \qquad \forall \beta,
\]
and therefore
\[
  \partial_t f(x,t)
  =
  -
  \E_{(x',y)}
  \Bigl[
    \frac{K_{\mathrm d}(x,x';t)+(M-1)K_{\mathrm o}(x,x';t)}{M}
    \partial_1 \ell(f(x',t),y)
  \Bigr].
\]
This is \eqref{eq:reduced}--\eqref{eq:keff}.
\end{proof}

\begin{corollary}[Equality of the two losses on the symmetric manifold]
\label{cor:losses}
Under the hypotheses of \cref{cor:collapse},
\[
  \riskhead(\rho_t)=\riskmean(\rho_t)
  \qquad
  \forall t \ge 0.
\]
\end{corollary}

\begin{proof}
By \cref{cor:collapse},
\[
  f_1(x;\rho_t)=\cdots=f_M(x;\rho_t)=f(x,t).
\]
Hence
\[
  \frac1M \sum_{\alpha=1}^M \ell(f_\alpha(x;\rho_t),y)=\ell(f(x,t),y)
\]
and also
\[
  \ell\Bigl(\frac1M \sum_{\alpha=1}^M f_\alpha(x;\rho_t),y\Bigr)=\ell(f(x,t),y).
\]
Taking expectations yields the claim.
\end{proof}

\begin{remark}
\Cref{cor:collapse,cor:losses} show that, at deterministic fixed-$M$ mean-field level,
symmetric initialization and a symmetric objective collapse headwise diversity completely.
This does \emph{not} imply that TabM is equivalent to a plain single MLP:
the effective kernel $K_{\mathrm{eff}}$ still depends on the shared/private parameterization.
It only says that the deterministic PDE no longer distinguishes the heads.
\end{remark}

\subsection{Nonlinear stability around the symmetric manifold}

\begin{definition}[Symmetrization]
For $\rho \in \PP_1(\Omega)$ define
\begin{equation}
  \Sym \rho
  =
  \frac1{M!} \sum_{\pi \in S_M} \Pi_{\pi\#}\rho.
  \label{eq:symmetrization}
\end{equation}
\end{definition}

\begin{theorem}[Nonlinear transverse stability of the symmetric manifold]
\label{thm:nonlinear-stability}
Assume \cref{ass:regularity}.
Let $\rho_t$ be the unique characteristic solution of \eqref{eq:pde} with initial law $\rho_0$,
and let $\rho_t^{\mathrm{sym}}$ be the unique characteristic solution with initial law
\[
  \rho_0^{\mathrm{sym}}=\Sym \rho_0.
\]
Then:
\begin{enumerate}[leftmargin=2em]
  \item $\rho_t^{\mathrm{sym}}$ is permutation-invariant for every $t\ge 0$;
  \item for every $t\ge 0$,
  \begin{equation}
    \W(\rho_t,\rho_t^{\mathrm{sym}})
    \le
    e^{\Lambda t}\W(\rho_0,\rho_0^{\mathrm{sym}});
    \label{eq:nonlinear-stability}
  \end{equation}
  \item if
  \[
    f^{\mathrm{sym}}(x,t)=f_1(x;\rho_t^{\mathrm{sym}}),
  \]
  then, for every $\alpha,\beta$ and every $x\in \mathcal X$,
  \begin{equation}
    |f_\alpha(x;\rho_t)-f^{\mathrm{sym}}(x,t)|
    \le
    L_0 e^{\Lambda t}\W(\rho_0,\rho_0^{\mathrm{sym}}),
    \label{eq:transverse-output}
  \end{equation}
  \begin{equation}
    |f_\alpha(x;\rho_t)-f_\beta(x;\rho_t)|
    \le
    2L_0 e^{\Lambda t}\W(\rho_0,\rho_0^{\mathrm{sym}}).
    \label{eq:transverse-disagreement}
  \end{equation}
\end{enumerate}
\end{theorem}

\begin{proof}
By construction, $\rho_0^{\mathrm{sym}}$ is invariant under every permutation.
Hence \cref{thm:symmetry} gives
\[
  \Pi_{\pi\#}\rho_t^{\mathrm{sym}}=\rho_t^{\mathrm{sym}}
  \qquad
  \forall \pi\in S_M,\ \forall t\ge 0,
\]
which proves the first claim.

The second claim is exactly \cref{prop:stability} applied to the two solutions
$\rho_t$ and $\rho_t^{\mathrm{sym}}$.

For the third claim, apply \cref{prop:lipschitz}:
\[
  |f_\alpha(x;\rho_t)-f_\alpha(x;\rho_t^{\mathrm{sym}})|
  \le
  L_0 \W(\rho_t,\rho_t^{\mathrm{sym}})
  \le
  L_0 e^{\Lambda t}\W(\rho_0,\rho_0^{\mathrm{sym}}).
\]
Because $\rho_t^{\mathrm{sym}}$ is permutation-invariant, \cref{cor:collapse} implies
\[
  f_\alpha(x;\rho_t^{\mathrm{sym}})
  =
  f^{\mathrm{sym}}(x,t)
  \qquad
  \forall \alpha.
\]
This proves \eqref{eq:transverse-output}.
Finally,
\[
  |f_\alpha(x;\rho_t)-f_\beta(x;\rho_t)|
  \le
  |f_\alpha(x;\rho_t)-f^{\mathrm{sym}}(x,t)|
  +
  |f_\beta(x;\rho_t)-f^{\mathrm{sym}}(x,t)|,
\]
and \eqref{eq:transverse-output} yields \eqref{eq:transverse-disagreement}.
\end{proof}

\begin{remark}
\Cref{thm:nonlinear-stability} is stronger than the exact-collapse theorem in one useful sense:
it shows that the symmetric manifold is not merely invariant,
but also controls nearby asymmetric initial laws quantitatively.
Therefore, in deterministic fixed-$M$ mean-field,
order-one headwise diversity cannot appear unless the initial law is itself order-one far from its symmetrization.
\end{remark}

\section{Architecture-Faithful Description and the Level-B Target}

\subsection{What BatchEnsemble and TabM actually implement}

To fix terminology, we briefly recall the architectural objects that the official papers and code implement.
In a BatchEnsemble layer, member $\alpha$ applies a shared weight matrix together with two member-specific modulation vectors:
\begin{equation}
  h_{\ell+1}^{(\alpha)}
  =
  \phi\Bigl(
    \bigl(
      W_\ell(h_\ell^{(\alpha)}\odot r_{\ell,\alpha})
    \bigr)
    \odot s_{\ell,\alpha}
    +
    b_\ell
  \Bigr).
  \label{eq:be-layer}
\end{equation}
Here $r_{\ell,\alpha}$ and $s_{\ell,\alpha}$ are vectors, not scalars,
and they act \emph{inside} the layer rather than as a single post-activation multiplier
\cite{batchensemble}.
The official TabM implementation uses a BatchEnsemble-style hidden backbone together with an independent output head for each member,
schematically
\begin{equation}
  y_\alpha
  =
  c_\alpha^\top h_L^{(\alpha)} + b_\alpha^{\mathrm{out}},
  \label{eq:tabm-arch-schematic}
\end{equation}
where the final readout parameters $(c_\alpha,b_\alpha^{\mathrm{out}})$ are not shared across $\alpha$
\cite{tabm,tabm-repo}.

\begin{remark}
Equation \eqref{eq:be-layer} is the architecture-faithful object.
The archived reduced-surrogate analyses in the companion note are \emph{not} exact rewritings of \eqref{eq:be-layer} or \eqref{eq:tabm-arch-schematic}.
They keep only three structural ingredients:
\begin{enumerate}[leftmargin=2em]
  \item one shared feature coordinate,
  \item member-specific feature modulation,
  \item independent versus shared output coordinates.
\end{enumerate}
Their purpose is to isolate coupling mechanisms in a form that fits cleanly into fixed-$M$ mean-field calculations.
\end{remark}

\subsection{Which model should carry the mean-field theory?}

The main modeling choice is not whether one wants a faithful architecture or a tractable theorem:
one ultimately wants both.
The real choice is which \emph{intermediate target} should carry the first rigorous mean-field analysis.
For the main text, it is enough to distinguish two levels.

\paragraph{Level A: architecture-faithful deep model.}
This is the layerwise system \eqref{eq:be-layer}--\eqref{eq:tabm-arch-schematic}.
It is the right long-term target, but it is also the most technically demanding one:
the shared weights, layerwise modulation vectors, and depth all interact.
A direct mean-field treatment would require a genuinely multilayer analysis rather than a single two-layer McKean--Vlasov equation.

\paragraph{Level B: semifaithful two-layer MLP-style targets.}
If one insists on a two-layer formulation while staying closer to the actual architecture, a convenient BatchEnsemble-like reduction is
\begin{equation}
  f_{\alpha}^{\mathrm{BE},m}(x)
  =
  \frac1m\sum_{j=1}^m
  c_j\,\sigma\bigl(w_j^\top (r_\alpha\odot x)+b_j\bigr),
  \label{eq:be-2layer}
\end{equation}
where $r_\alpha\in\R^d$ is a head-specific input modulation vector shared across all hidden units.
In two layers, any final hidden-side modulation can be absorbed into the readout, so \eqref{eq:be-2layer} should be read as a \emph{reduced} two-layer proxy rather than as a literal transcription of \eqref{eq:be-layer}.
A natural TabM-like variant is
\begin{equation}
  f_{\alpha}^{\mathrm{TabM},m}(x)
  =
  \frac1m\sum_{j=1}^m
  a_{\alpha,j}\,\sigma\bigl(w_j^\top (r_\alpha\odot x)+b_j\bigr).
  \label{eq:tabm-2layer}
\end{equation}
These are much closer to the real architecture than the archived companion-note surrogates.
However, from the mean-field viewpoint they already introduce an extra complication:
the vectors $r_\alpha$ are \emph{global head parameters}, not per-neuron coordinates.
So the natural limit is no longer just an empirical measure on neuron states;
it is a coupled system consisting of a neuron measure plus finitely many global head variables.

\begin{remark}
For the purposes of the present note, Level B is the main target.
It is already more complicated than the standard two-layer theory because one has to track both
\begin{enumerate}[leftmargin=2em]
  \item the empirical measure of hidden neurons, and
  \item the finite-dimensional global head modulations $(r_\alpha)_{\alpha=1}^M$.
\end{enumerate}
Level A is harder still, because depth prevents the entire dynamics from closing at the level of a single two-layer particle law.
So the instinct that ``mean-field is safest in two layers'' is essentially correct, but within two layers the semifaithful Level-B target remains the most relevant one.
A separate companion note records more aggressive Level-C reduced-surrogate calculations that are omitted from the main text to keep the story architecture-faithful.
\end{remark}

\subsection{A Level-B coupled mean-field system}

The main semifaithful two-layer TabM-like target in this note is \eqref{eq:tabm-2layer}.
Write the neuron state as
\[
  \vartheta=(w,b,(a_\alpha)_{\alpha=1}^M)
  \in
  \Theta
  :=
  \R^d\times \R \times \R^M,
\]
and the global head modulations as
\[
  r=(r_\alpha)_{\alpha=1}^M\in \mathcal{R}:=(\R^d)^M.
\]
For a neuron law $\rho\in \PP_1(\Theta)$ and a head-modulation vector $r\in\mathcal R$, define
\[
  \Psi_\alpha(x;\vartheta,r)
  :=
  a_\alpha \sigma\bigl(w^\top (r_\alpha\odot x)+b\bigr),
\]
so that
\begin{equation}
  f_\alpha(x;\rho,r)
  :=
  \int_\Theta
  \Psi_\alpha(x;\vartheta,r)
  \,\rho(\dd \vartheta).
  \label{eq:levelb-output}
\end{equation}
The corresponding per-head risk is
\begin{equation}
  \mathcal R_{\mathrm{LevelB}}(\rho,r)
  :=
  \E_{(x,y)\sim \mathsf P}
  \Bigl[
    \frac1M\sum_{\alpha=1}^M
    \ell(f_\alpha(x;\rho,r),y)
  \Bigr].
  \label{eq:levelb-risk}
\end{equation}

\paragraph{Finite-width motivation.}
For width $m$, consider
\[
  \rho_t^m=\frac1m\sum_{j=1}^m \delta_{\vartheta_j(t)},
\]
where
\[
  f_{\alpha}^{m}(x)
  =
  \frac1m\sum_{j=1}^m
  a_{\alpha,j}\,\sigma\bigl(w_j^\top (r_\alpha\odot x)+b_j\bigr).
\]
Under gradient flow for the population risk, the neuron gradients scale like $m^{-1}$ while the head-global gradients scale like $1$.
So the natural coupled scaling is
\begin{equation}
  \dot \vartheta_j = -m \nabla_{\vartheta_j}\mathcal R_m,
  \qquad
  \dot r_\alpha = - \nabla_{r_\alpha}\mathcal R_m.
  \label{eq:levelb-scaling}
\end{equation}
This is the first conceptual difference from the standard particle-only two-layer mean-field setup:
the mean-field limit is no longer a pure PDE on a neuron measure, but a PDE for $\rho_t$ coupled to an ODE for $r_t$.

\begin{definition}[Level-B particle and head drifts]
For $(\rho,r)\in \PP_1(\Theta)\times \mathcal R$, define the neuron potential
\begin{equation}
  V(\vartheta;\rho,r)
  :=
  \E_{(x,y)\sim\mathsf P}
  \Bigl[
    \frac1M\sum_{\alpha=1}^M
    \partial_1\ell(f_\alpha(x;\rho,r),y)\,
    a_\alpha \sigma\bigl(w^\top(r_\alpha\odot x)+b\bigr)
  \Bigr],
  \label{eq:levelb-potential}
\end{equation}
the neuron drift
\begin{equation}
  B(\vartheta;\rho,r):=-\nabla_\vartheta V(\vartheta;\rho,r),
  \label{eq:levelb-neuron-drift}
\end{equation}
and the head drift
\begin{equation}
  G_\alpha(\rho,r)
  :=
  -\E_{(x,y)\sim\mathsf P}
  \Bigl[
    \frac1M
    \partial_1\ell(f_\alpha(x;\rho,r),y)
    \int_\Theta
    a_\alpha \sigma'\bigl(w^\top(r_\alpha\odot x)+b\bigr)
    (w\odot x)
    \,\rho(\dd \vartheta)
  \Bigr].
  \label{eq:levelb-head-drift}
\end{equation}
\end{definition}

The formal Level-B mean-field dynamics is therefore
\begin{equation}
  \partial_t \rho_t
  +
  \nabla_\vartheta\cdot\bigl(\rho_t B(\vartheta;\rho_t,r_t)\bigr)
  =
  0,
  \label{eq:levelb-pde}
\end{equation}
\begin{equation}
  \dot r_{\alpha,t}=G_\alpha(\rho_t,r_t),
  \qquad
  \alpha=1,\dots,M.
  \label{eq:levelb-ode}
\end{equation}

\begin{theorem}[Truncated Level-B characteristic well-posedness and symmetry]
\label{thm:levelb-wellposed}
Assume that the input domain is bounded:
\[
  \|x\|\le X_*
  \qquad
  \forall x\in \mathcal X,
\]
that $\sigma\in C_b^2(\R)$, and that $\partial_1\ell(\cdot,y)$ is globally Lipschitz and uniformly bounded in $y$.
Fix radii $R_w,R_b,R_a,R_r>0$ and define the compact identity region
\[
  \Theta_R
  :=
  \bigl\{
    (w,b,(a_\alpha)_\alpha):
    \|w\|\le R_w,\ |b|\le R_b,\ |a_\alpha|\le R_a
  \bigr\},
\]
\[
  \mathcal R_R
  :=
  \bigl\{
    (r_\alpha)_\alpha:
    \|r_\alpha\|\le R_r
    \ \forall \alpha
  \bigr\}.
\]
Choose smooth truncation maps
\[
  T_{\Theta,R}:\Theta\to\Theta,
  \qquad
  T_{\mathcal R,R}:\mathcal R\to\mathcal R,
\]
with bounded images and globally bounded and globally Lipschitz first derivatives, such that
\[
  T_{\Theta,R}(\vartheta)=\vartheta
  \quad
  \forall \vartheta\in \Theta_R,
  \qquad
  T_{\mathcal R,R}(r)=r
  \quad
  \forall r\in \mathcal R_R.
\]
Assume in addition that the truncation maps commute with head permutations:
for every $\pi\in S_M$,
\[
  T_{\Theta,R}(\Pi_\pi \vartheta)=\Pi_\pi T_{\Theta,R}(\vartheta),
  \qquad
  T_{\mathcal R,R}(\Pi_\pi r)=\Pi_\pi T_{\mathcal R,R}(r).
\]
Write
\[
  T_{\Theta,R}(\vartheta)=(\widetilde w,\widetilde b,(\widetilde a_\alpha)_\alpha),
  \qquad
  T_{\mathcal R,R}(r)=(\widetilde r_\alpha)_\alpha,
\]
and define the truncated feature
\[
  \Psi_{\alpha,R}(x;\vartheta,r)
  :=
  \widetilde a_\alpha
  \sigma\bigl(\widetilde w^\top(\widetilde r_\alpha\odot x)+\widetilde b\bigr),
\]
and denote the corresponding truncated outputs, drifts, and coupled dynamics by
\[
  f_{\alpha,R},\qquad B_R,\qquad G_{\alpha,R},
\]
and
\[
  G_R:=(G_{\alpha,R})_{\alpha=1}^M\in \mathcal R.
\]
That is, one replaces $\Psi_\alpha$ with $\Psi_{\alpha,R}$ in \eqref{eq:levelb-output}--\eqref{eq:levelb-ode}.
Then the truncated coupled system has a unique global characteristic solution
\[
  (\rho_t,r_t)_{t\ge 0}
  \in
  C\bigl([0,\infty),\PP_1(\Theta)\times \mathcal R\bigr)
\]
for every initial condition $(\rho_0,r_0)\in \PP_1(\Theta)\times \mathcal R$.
More precisely, there exists a unique flow map $\Phi_t:\Theta\to\Theta$ such that
\[
  \dot \Phi_t(\vartheta)=B_R(\Phi_t(\vartheta);\rho_t,r_t),
  \qquad
  \Phi_0(\vartheta)=\vartheta,
  \qquad
  \rho_t=(\Phi_t)_\#\rho_0,
\]
and
\[
  \dot r_t=G_R(\rho_t,r_t).
\]

Moreover, if $\rho_0$ is exchangeable across the coordinates $(a_\alpha)_{\alpha=1}^M$ and
\[
  r_{1,0}=\cdots=r_{M,0},
\]
then the solution remains permutation-symmetric and
\[
  f_{1,R}(x;\rho_t,r_t)=\cdots=f_{M,R}(x;\rho_t,r_t)
  \qquad
  \forall x\in \mathcal X,\ \forall t\ge 0.
\]
Finally, if for some $T>0$ the truncated solution satisfies
\[
  \supp(\rho_t)\subseteq \Theta_R,
  \qquad
  r_t\in \mathcal R_R
  \qquad
  \forall t\in[0,T],
\]
then on $[0,T]$ it also solves the original untruncated Level-B system \eqref{eq:levelb-pde}--\eqref{eq:levelb-ode}.
\end{theorem}

\begin{proof}
The smooth truncation maps have bounded images and globally bounded and globally Lipschitz first derivatives, so the truncated feature $\Psi_{\alpha,R}$ and its first derivatives with respect to $\vartheta$ and $r_\alpha$ are globally bounded and globally Lipschitz.
Therefore the truncated outputs are Lipschitz in $(\rho,r)$ with respect to the metric
\[
  \mathbf d\bigl((\rho,r),(\mu,s)\bigr)
  :=
  \W(\rho,\mu)+\|r-s\|,
\]
and the corresponding truncated drifts are jointly Lipschitz in $(\vartheta,\rho,r)$ and $(\rho,r)$, respectively.

Given a continuous path $(\widehat\rho_t,\widehat r_t)$, solve
\[
  \dot \Theta_t = B_R(\Theta_t;\widehat\rho_t,\widehat r_t),
  \qquad
  \Theta_0\sim \rho_0,
\]
and
\[
  \dot r_{\alpha,t} = G_{\alpha,R}(\widehat\rho_t,\widehat r_t),
  \qquad
  r_{\alpha,0}=(r_0)_\alpha,
  \qquad
  \alpha=1,\dots,M.
\]
This defines a map on path space by
\[
  (\widehat\rho_t,\widehat r_t)
  \mapsto
  (\Theta_{t\#}\rho_0,r_t).
\]
The Lipschitz bounds above imply, exactly as in the proof of \cref{thm:wellposed}, that this map is a contraction on a short time interval in the metric
\[
  \sup_{t\in[0,T]}
  \mathbf d\bigl((\rho_t,r_t),(\mu_t,s_t)\bigr).
\]
Iterating the local contraction yields a unique global characteristic solution together with its characteristic flow $\Phi_t$.

For symmetry, let $\pi\in S_M$ act by permuting the coordinates $(a_\alpha)_\alpha$ of $\vartheta$ and the head vectors $(r_\alpha)_\alpha$.
Because
\[
  \Psi_{\alpha,R}(x;\Pi_\pi \vartheta,\Pi_\pi r)
  =
  \Psi_{\pi^{-1}(\alpha),R}(x;\vartheta,r),
\]
the coupled drifts are equivariant under this joint permutation;
here we use both the permutation-equivariance of the model itself and the commutation of the truncation maps with $\Pi_\pi$.
Hence if $\rho_0$ is exchangeable and $r_{1,0}=\cdots=r_{M,0}$, then the solution remains invariant under every permutation.
The outputs must therefore coincide across $\alpha$.

Finally, if the truncated solution remains inside $\Theta_R\times \mathcal R_R$ on $[0,T]$, then $T_{\Theta,R}$ and $T_{\mathcal R,R}$ act as the identity along the trajectory.
Hence the truncated and untruncated drifts coincide there, so the same pair $(\rho_t,r_t)$ solves the original Level-B system on $[0,T]$ as well.
\end{proof}

\begin{proposition}[Stability of the truncated Level-B flow]
\label{prop:levelb-stability}
Under the assumptions of \cref{thm:levelb-wellposed}, there exists a constant $\Lambda_R<\infty$,
depending only on the truncation radii and the regularity constants, such that the following holds.
Let $(\rho_t,r_t)$ and $(\mu_t,s_t)$ be the unique solutions of the truncated Level-B system
with initial data $(\rho_0,r_0)$ and $(\mu_0,s_0)$, respectively.
Then, for every $t\ge 0$,
\begin{equation}
  \mathbf d\bigl((\rho_t,r_t),(\mu_t,s_t)\bigr)
  \le
  e^{\Lambda_R t}
  \mathbf d\bigl((\rho_0,r_0),(\mu_0,s_0)\bigr),
  \label{eq:levelb-stability}
\end{equation}
where
\[
  \mathbf d\bigl((\rho,r),(\mu,s)\bigr)
  :=
  \W(\rho,\mu)+\|r-s\|.
\]
\end{proposition}

\begin{proof}
By the proof of \cref{thm:levelb-wellposed}, the truncated drifts satisfy global Lipschitz bounds of the form
\begin{equation}
  \|B_R(\vartheta;\rho,r)-B_R(\eta;\mu,s)\|
  \le
  L_{B,R}
  \bigl(
    \|\vartheta-\eta\|
    +
    \W(\rho,\mu)
    +
    \|r-s\|
  \bigr),
  \label{eq:levelb-br-lip}
\end{equation}
and
\begin{equation}
  \|G_R(\rho,r)-G_R(\mu,s)\|
  \le
  L_{G,R}
  \bigl(
    \W(\rho,\mu)
    +
    \|r-s\|
  \bigr),
  \label{eq:levelb-gr-lip}
\end{equation}
where
\[
  G_R(\rho,r):=(G_{\alpha,R}(\rho,r))_{\alpha=1}^M\in \mathcal R.
\]

Let $\Phi_t$ and $\Psi_t$ be the characteristic flows generated by
$B_R(\cdot;\rho_t,r_t)$ and $B_R(\cdot;\mu_t,s_t)$, respectively.
Fix any coupling $\gamma_0$ of $\rho_0$ and $\mu_0$, and define
\[
  \gamma_t=(\Phi_t,\Psi_t)_\#\gamma_0.
\]
Then $\gamma_t$ is a coupling of $\rho_t$ and $\mu_t$, so
\[
  \W(\rho_t,\mu_t)
  \le
  D_t^\vartheta
  :=
  \int_{\Theta\times\Theta}
  \|\Phi_t(\vartheta)-\Psi_t(\eta)\|
  \,\gamma_0(\dd \vartheta,\dd \eta).
\]
Set
\[
  D_t:=D_t^\vartheta+\|r_t-s_t\|.
\]
Differentiating $D_t^\vartheta$ and using \eqref{eq:levelb-br-lip} gives
\begin{align*}
  \frac{\dd}{\dd t}D_t^\vartheta
  &\le
  \int_{\Theta\times\Theta}
  \|B_R(\Phi_t(\vartheta);\rho_t,r_t)-B_R(\Psi_t(\eta);\mu_t,s_t)\|
  \,\gamma_0(\dd \vartheta,\dd \eta)
  \\
  &\le
  L_{B,R}
  \Bigl(
    D_t^\vartheta+\W(\rho_t,\mu_t)+\|r_t-s_t\|
  \Bigr)
  \\
  &\le
  2L_{B,R}D_t^\vartheta+L_{B,R}\|r_t-s_t\|,
\end{align*}
where we used $\W(\rho_t,\mu_t)\le D_t^\vartheta$.
Similarly, \eqref{eq:levelb-gr-lip} yields
\[
  \frac{\dd}{\dd t}\|r_t-s_t\|
  \le
  L_{G,R}\bigl(\W(\rho_t,\mu_t)+\|r_t-s_t\|\bigr)
  \le
  L_{G,R}D_t.
\]
Therefore
\[
  \dot D_t
  \le
  (2L_{B,R}+L_{G,R})D_t^\vartheta+(L_{B,R}+L_{G,R})\|r_t-s_t\|
  \le
  (2L_{B,R}+L_{G,R})D_t.
\]
Hence, with $\Lambda_R:=2L_{B,R}+L_{G,R}$, Gronwall's inequality gives
\[
  D_t\le e^{\Lambda_R t}D_0.
\]
Since
\[
  \mathbf d\bigl((\rho_t,r_t),(\mu_t,s_t)\bigr)
  \le D_t,
\]
taking the infimum over all initial couplings $\gamma_0$ proves \eqref{eq:levelb-stability}.
\end{proof}

\begin{corollary}[Deterministic finite-width convergence for the truncated Level-B system]
\label{cor:levelb-empirical-limit}
Under the assumptions of \cref{thm:levelb-wellposed}, let
\[
  \rho_0^m=\frac1m\sum_{j=1}^m \delta_{\vartheta_j^m(0)}
  \in \PP_1(\Theta),
  \qquad
  r_0^m\in \mathcal R,
\]
and consider the finite-width truncated particle/head system
\begin{equation}
  \dot \vartheta_j^m(t)
  =
  B_R(\vartheta_j^m(t);\rho_t^m,r_t^m),
  \qquad
  \rho_t^m=\frac1m\sum_{j=1}^m \delta_{\vartheta_j^m(t)},
  \label{eq:levelb-particle-system}
\end{equation}
\begin{equation}
  \dot r_t^m
  =
  G_R(\rho_t^m,r_t^m),
  \qquad
  r_t^m\in \mathcal R.
  \label{eq:levelb-head-system}
\end{equation}
Let $(\rho_t,r_t)$ be the unique truncated Level-B characteristic solution with initial data $(\rho_0,r_0)$.
Then, for every $T>0$,
\begin{equation}
  \sup_{t\in[0,T]}
  \mathbf d\bigl((\rho_t^m,r_t^m),(\rho_t,r_t)\bigr)
  \le
  e^{\Lambda_R T}
  \mathbf d\bigl((\rho_0^m,r_0^m),(\rho_0,r_0)\bigr),
  \label{eq:levelb-empirical}
\end{equation}
where $\Lambda_R$ is the constant from \cref{prop:levelb-stability}.
\end{corollary}

\begin{proof}
The pair $(\rho_t^m,r_t^m)$ is the pushforward solution generated by the truncated drift $B_R(\cdot;\rho_t^m,r_t^m)$ together with the head ODE \eqref{eq:levelb-head-system},
so it is a solution of the truncated Level-B coupled system with initial data $(\rho_0^m,r_0^m)$.
By uniqueness in \cref{thm:levelb-wellposed}, it is exactly the unique truncated characteristic solution with that initial condition.
Applying \cref{prop:levelb-stability} to the pair
\[
  (\rho_t^m,r_t^m)
  \qquad \text{and} \qquad
  (\rho_t,r_t)
\]
gives \eqref{eq:levelb-empirical}.
\end{proof}

\begin{corollary}[Headwise disagreement vanishes in the truncated Level-B deterministic regime]
\label{cor:levelb-disagreement}
Under the assumptions of \cref{thm:levelb-wellposed}, let $(\rho_t,r_t)$ be the truncated Level-B solution with symmetric initial data:
$\rho_0$ exchangeable across $(a_\alpha)_{\alpha=1}^M$ and
\[
  r_{1,0}=\cdots=r_{M,0}.
\]
Let $(\rho_t^m,r_t^m)$ be the finite-width truncated system from \cref{cor:levelb-empirical-limit}.
Then there exists a constant $L_{\Psi,R}<\infty$ such that, for every $T>0$,
\begin{equation}
  \sup_{t\in[0,T]}
  \sup_{x\in\mathcal X}
  \max_{\alpha,\beta\le M}
  |f_{\alpha,R}(x;\rho_t^m,r_t^m)-f_{\beta,R}(x;\rho_t^m,r_t^m)|
  \le
  2L_{\Psi,R}e^{\Lambda_R T}
  \mathbf d\bigl((\rho_0^m,r_0^m),(\rho_0,r_0)\bigr).
  \label{eq:levelb-disagreement}
\end{equation}
In particular, if
\[
  \mathbf d\bigl((\rho_0^m,r_0^m),(\rho_0,r_0)\bigr)\to 0,
\]
then deterministic headwise disagreement vanishes uniformly on every compact time interval.
\end{corollary}

\begin{proof}
Because $\Psi_{\alpha,R}(x;\vartheta,r)$ is globally Lipschitz in $(\vartheta,r)$ uniformly over $x\in\mathcal X$,
there exists $L_{\Psi,R}<\infty$ such that
\[
  |f_{\alpha,R}(x;\rho,r)-f_{\alpha,R}(x;\mu,s)|
  \le
  L_{\Psi,R}\mathbf d\bigl((\rho,r),(\mu,s)\bigr)
\]
for all $x,\alpha$ and all $(\rho,r),(\mu,s)$.
By the symmetry part of \cref{thm:levelb-wellposed},
\[
  f_{1,R}(x;\rho_t,r_t)=\cdots=f_{M,R}(x;\rho_t,r_t)
  \qquad \forall x,\ t.
\]
Hence, for every $\alpha,\beta$,
\begin{align*}
  |f_{\alpha,R}(x;\rho_t^m,r_t^m)-f_{\beta,R}(x;\rho_t^m,r_t^m)|
  &\le
  |f_{\alpha,R}(x;\rho_t^m,r_t^m)-f_{\alpha,R}(x;\rho_t,r_t)|
  \\
  &\quad+
  |f_{\beta,R}(x;\rho_t^m,r_t^m)-f_{\beta,R}(x;\rho_t,r_t)|
  \\
  &\le
  2L_{\Psi,R}\mathbf d\bigl((\rho_t^m,r_t^m),(\rho_t,r_t)\bigr).
\end{align*}
Now apply \cref{cor:levelb-empirical-limit}.
\end{proof}

\begin{corollary}[Initialization rate implies a width law for deterministic disagreement]
\label{cor:levelb-width-law}
Under the assumptions of \cref{cor:levelb-disagreement}, suppose that the initial discrepancy satisfies
\[
  \mathbf d\bigl((\rho_0^m,r_0^m),(\rho_0,r_0)\bigr)\le \varepsilon_m
\]
for some deterministic sequence $\varepsilon_m\downarrow 0$.
Then, for every $T>0$,
\begin{equation}
  \sup_{t\in[0,T]}
  \sup_{x\in\mathcal X}
  \max_{\alpha,\beta\le M}
  |f_{\alpha,R}(x;\rho_t^m,r_t^m)-f_{\beta,R}(x;\rho_t^m,r_t^m)|
  \le
  2L_{\Psi,R}e^{\Lambda_R T}\,\varepsilon_m.
  \label{eq:levelb-width-law}
\end{equation}
In particular, any CLT-scale initialization error
\[
  \varepsilon_m=O(m^{-1/2})
\]
implies CLT-scale deterministic head disagreement on every compact time interval.
\end{corollary}

\begin{proof}
This is exactly \cref{eq:levelb-disagreement} with the assumed bound on the initial discrepancy.
\end{proof}

\begin{remark}
\Cref{thm:levelb-wellposed,prop:levelb-stability,cor:levelb-empirical-limit,cor:levelb-disagreement}
show that the semifaithful Level-B model is already mathematically tractable as a \emph{truncated} deterministic mean-field theory.
The present proof uses truncation only to force global bounded-Lipschitz estimates in the metric
\[
  \mathbf d\bigl((\rho,r),(\mu,s)\bigr)=\W(\rho,\mu)+\|r-s\|.
\]
For the untruncated Level-B system, the genuine obstruction is the head drift \cref{eq:levelb-head-drift},
which contains the mixed observable $a_\alpha w$ and therefore is not globally bounded-Lipschitz in that metric.
This suggests that truncation removal should be treated as a finite-horizon moment problem rather than as a new modeling problem.
More precisely, if the initial law has essentially bounded head coordinates $(a_\alpha)_\alpha$ and finite second moments in $(w,b)$,
then one expects uniform bounds on $(a_\alpha(t))_\alpha$, a coupled Gronwall estimate for
\[
  \int_\Theta \|w\|\,\rho_t(\dd \vartheta)
  \qquad \text{and} \qquad
  \|r_t\|,
\]
propagation of second moments, and hence a local contraction on bounded-moment sets in a metric such as $W_2+\|\cdot\|$.
That finite-horizon untruncated upgrade is not proved in this note,
but it now looks like a technically plausible next theorem rather than a conceptual obstruction.
\end{remark}

\begin{definition}[Regular characteristic solutions for untruncated Level B]
\label{def:levelb-regular-characteristic}
Let $(\rho_t,r_t)_{t\in[0,T]}$ solve the untruncated Level-B system.
We call it a \emph{regular characteristic solution} on $[0,T]$ if there exists a characteristic flow
\[
  \Phi:[0,T]\times\Theta\to\Theta
\]
such that
\[
  \dot \Phi_t(\vartheta)=B(\Phi_t(\vartheta);\rho_t,r_t),
  \qquad
  \Phi_0(\vartheta)=\vartheta,
  \qquad
  \rho_t=(\Phi_t)_\#\rho_0,
\]
the maps $t\mapsto \Phi_t(\vartheta)$ and $t\mapsto r_t$ are $C^1$ for $\rho_0$-almost every $\vartheta$,
and the observables appearing below are integrable uniformly enough to justify differentiation under the integral sign in \eqref{eq:levelb-output}.
The same terminology is used for the shared-output proxy after replacing the corresponding feature maps and drifts in the obvious way.
\end{definition}

\begin{proposition}[Level-B output dynamics has a neuron kernel and a head kernel]
\label{prop:levelb-dynamics}
Assume that $(\rho_t,r_t)_{t\in[0,T]}$ is a regular characteristic solution of the untruncated Level-B system \eqref{eq:levelb-pde}--\eqref{eq:levelb-ode} in the sense of \cref{def:levelb-regular-characteristic}. Define
\[
  z_\alpha(x;\vartheta,r):=w^\top(r_\alpha\odot x)+b,
  \qquad
  g_\alpha(x;\rho,r)
  :=
  \int_\Theta
  a_\alpha \sigma'(z_\alpha(x;\vartheta,r))(w\odot x)
  \,\rho(\dd \vartheta).
\]
Then the outputs satisfy
\begin{align}
  \partial_t f_\alpha(x;\rho_t,r_t)
  &=
  -\E_{(x',y)\sim\mathsf P}
  \Bigl[
    \frac1M\sum_{\beta=1}^M
    K^{\mathrm{neu}}_{\alpha\beta}(x,x';\rho_t,r_t)
    \partial_1\ell(f_\beta(x';\rho_t,r_t),y)
  \Bigr]
  \nonumber\\
  &\quad
  -
  \E_{(x',y)\sim\mathsf P}
  \Bigl[
    \frac1M
    K^{\mathrm{head}}_{\alpha\alpha}(x,x';\rho_t,r_t)
    \partial_1\ell(f_\alpha(x';\rho_t,r_t),y)
  \Bigr],
  \label{eq:levelb-output-dynamics}
\end{align}
where
\begin{align}
  K^{\mathrm{neu}}_{\alpha\beta}(x,x';\rho,r)
  &=
  \int_\Theta
  a_\alpha a_\beta
  \sigma'(z_\alpha(x;\vartheta,r))
  \sigma'(z_\beta(x';\vartheta,r))
  \Bigl[
    1+(r_\alpha\odot x)^\top (r_\beta\odot x')
  \Bigr]
  \rho(\dd \vartheta)
  \nonumber\\
  &\quad
  +
  \delta_{\alpha\beta}
  \int_\Theta
  \sigma(z_\alpha(x;\vartheta,r))
  \sigma(z_\alpha(x';\vartheta,r))
  \rho(\dd \vartheta),
  \label{eq:levelb-neuron-kernel}
\end{align}
and
\begin{equation}
  K^{\mathrm{head}}_{\alpha\alpha}(x,x';\rho,r)
  :=
  \langle g_\alpha(x;\rho,r),g_\alpha(x';\rho,r)\rangle_{\R^d}.
  \label{eq:levelb-head-kernel}
\end{equation}
\end{proposition}

\begin{proof}
Differentiate \eqref{eq:levelb-output} along the coupled dynamics:
\[
  \partial_t f_\alpha(x;\rho_t,r_t)
  =
  \int_\Theta
  \langle \nabla_\vartheta \Psi_\alpha(x;\vartheta,r_t),
  B(\vartheta;\rho_t,r_t)\rangle
  \rho_t(\dd \vartheta)
  +
  \bigl\langle \nabla_{r_\alpha}f_\alpha(x;\rho_t,r_t),\dot r_{\alpha,t}\bigr\rangle.
\]
Because
\[
  \partial_{a_\gamma}\Psi_\alpha
  =
  \delta_{\alpha\gamma}\sigma(z_\alpha),
  \qquad
  \partial_b\Psi_\alpha
  =
  a_\alpha \sigma'(z_\alpha),
  \qquad
  \nabla_w\Psi_\alpha
  =
  a_\alpha \sigma'(z_\alpha)(r_\alpha\odot x),
\]
the neuron contribution expands into the kernel term \eqref{eq:levelb-neuron-kernel}.
The derivative with respect to the global head variable is
\[
  \nabla_{r_\alpha}f_\alpha(x;\rho,r)=g_\alpha(x;\rho,r).
\]
Substituting \eqref{eq:levelb-ode} gives
\[
  \bigl\langle \nabla_{r_\alpha}f_\alpha(x;\rho_t,r_t),\dot r_{\alpha,t}\bigr\rangle
  =
  -\E_{(x',y)\sim\mathsf P}
  \Bigl[
    \frac1M
    \langle g_\alpha(x;\rho_t,r_t),g_\alpha(x';\rho_t,r_t)\rangle
    \partial_1\ell(f_\alpha(x';\rho_t,r_t),y)
  \Bigr],
 \]
which is exactly the head-kernel term \eqref{eq:levelb-head-kernel}.
Combining the two pieces proves \eqref{eq:levelb-output-dynamics}.
\end{proof}

\begin{remark}
\Cref{prop:levelb-dynamics} makes the main structural point explicit.
At Level B, there is an additional \emph{head kernel} coming from the global modulation variables $(r_\alpha)_\alpha$.
This is precisely why Level B is both more faithful and more difficult.
\end{remark}

\begin{corollary}[Collapsed Level-B dynamics carries an extra head-kernel term]
\label{cor:levelb-collapsed}
Assume that $(\rho_t,r_t)_{t\in[0,T]}$ is a regular characteristic solution of the untruncated Level-B system such that, for every $t\in[0,T]$,
\[
  r_{1,t}=\cdots=r_{M,t}=:r_t^\star,
\]
and $\rho_t$ is exchangeable across the coordinates $(a_\alpha)_{\alpha=1}^M$.
Then
\[
  f_1(x;\rho_t,r_t)=\cdots=f_M(x;\rho_t,r_t)=:f(x,t),
  \qquad
  \forall x\in\mathcal X,
\]
and the common predictor evolves as
\begin{equation}
  \partial_t f(x,t)
  =
  -\E_{(x',y)\sim\mathsf P}
  \Bigl[
    K_{\mathrm{eff}}^{\mathrm{LevelB}}(x,x';\rho_t,r_t^\star)
    \partial_1\ell(f(x',t),y)
  \Bigr],
  \label{eq:levelb-keff-dynamics}
\end{equation}
with
\begin{equation}
  K_{\mathrm{eff}}^{\mathrm{LevelB}}
  :=
  \frac{
    K_d^{\mathrm{neu}}
    +(M-1)K_o^{\mathrm{neu}}
    +K^{\mathrm{head}}
  }{M},
  \label{eq:levelb-keff}
\end{equation}
where
\[
  K_d^{\mathrm{neu}}
  :=
  K_{\alpha\alpha}^{\mathrm{neu}},
  \qquad
  K_o^{\mathrm{neu}}
  :=
  K_{\alpha\beta}^{\mathrm{neu}}
  \ \ (\alpha\neq\beta),
  \qquad
  K^{\mathrm{head}}
  :=
  K_{\alpha\alpha}^{\mathrm{head}},
\]
which are independent of the particular choice of $\alpha,\beta$ by exchangeability.
\end{corollary}

\begin{proof}
The equality of the outputs follows from exchangeability of $\rho_t$ and the fact that all head modulations are equal to $r_t^\star$.
Under the same symmetry, the diagonal and off-diagonal neuron kernels are independent of the index choice, and so is the diagonal head kernel.
Substituting these identities into \eqref{eq:levelb-output-dynamics} and using $f_\beta=f$ for all $\beta$ yields \eqref{eq:levelb-keff-dynamics} with the effective kernel \eqref{eq:levelb-keff}.
\end{proof}

\subsection{Comparison with a shared-output Level-B proxy}

To isolate the effect of the independent output head, it is useful to compare the TabM-like Level-B model with the corresponding shared-output two-layer proxy
\begin{equation}
  f_{\alpha}^{\mathrm{sh},m}(x)
  =
  \frac1m\sum_{j=1}^m
  c_j\,\sigma\bigl(w_j^\top (r_\alpha\odot x)+b_j\bigr),
  \label{eq:levelb-shared-finite}
\end{equation}
which is the reduced two-layer BatchEnsemble-like model from \eqref{eq:be-2layer}.
Its neuron state is now
\[
  \vartheta=(w,b,c)\in \Theta^{\mathrm{sh}}:=\R^d\times \R\times \R.
\]
For $\rho\in \PP_1(\Theta^{\mathrm{sh}})$ and $r\in \mathcal R$, define
\[
  \Psi_\alpha^{\mathrm{sh}}(x;\vartheta,r)
  :=
  c\,\sigma\bigl(w^\top(r_\alpha\odot x)+b\bigr),
\]
\[
  f_\alpha^{\mathrm{sh}}(x;\rho,r)
  :=
  \int_{\Theta^{\mathrm{sh}}}
  \Psi_\alpha^{\mathrm{sh}}(x;\vartheta,r)\,\rho(\dd \vartheta),
\]
and the corresponding population risk
\[
  \mathcal R_{\mathrm{LevelB}}^{\mathrm{sh}}(\rho,r)
  :=
  \E_{(x,y)\sim\mathsf P}
  \Bigl[
    \frac1M\sum_{\alpha=1}^M
    \ell(f_\alpha^{\mathrm{sh}}(x;\rho,r),y)
  \Bigr].
\]

\begin{proposition}[Shared-output Level-B dynamics]
\label{prop:levelb-shared-dynamics}
Assume that $(\rho_t,r_t)_{t\in[0,T]}$ is a regular characteristic solution of the untruncated shared-output Level-B system associated with \eqref{eq:levelb-shared-finite}.
Define
\[
  z_\alpha(x;\vartheta,r):=w^\top(r_\alpha\odot x)+b,
  \qquad
  g_\alpha^{\mathrm{sh}}(x;\rho,r)
  :=
  \int_{\Theta^{\mathrm{sh}}}
  c\,\sigma'(z_\alpha(x;\vartheta,r))(w\odot x)
  \,\rho(\dd \vartheta).
\]
Then the outputs satisfy
\begin{align}
  \partial_t f_\alpha^{\mathrm{sh}}(x;\rho_t,r_t)
  &=
  -\E_{(x',y)\sim\mathsf P}
  \Bigl[
    \frac1M\sum_{\beta=1}^M
    K_{\alpha\beta}^{\mathrm{sh}}(x,x';\rho_t,r_t)
    \partial_1\ell(f_\beta^{\mathrm{sh}}(x';\rho_t,r_t),y)
  \Bigr]
  \nonumber\\
  &\quad
  -
  \E_{(x',y)\sim\mathsf P}
  \Bigl[
    \frac1M
    K_{\alpha\alpha}^{\mathrm{sh,head}}(x,x';\rho_t,r_t)
    \partial_1\ell(f_\alpha^{\mathrm{sh}}(x';\rho_t,r_t),y)
  \Bigr],
  \label{eq:levelb-shared-output-dynamics}
\end{align}
where
\begin{align}
  K_{\alpha\beta}^{\mathrm{sh}}(x,x';\rho,r)
  &=
  \int_{\Theta^{\mathrm{sh}}}
  c^2
  \sigma'(z_\alpha(x;\vartheta,r))
  \sigma'(z_\beta(x';\vartheta,r))
  \Bigl[
    1+(r_\alpha\odot x)^\top(r_\beta\odot x')
  \Bigr]
  \rho(\dd \vartheta)
  \nonumber\\
  &\quad
  +
  \int_{\Theta^{\mathrm{sh}}}
  \sigma(z_\alpha(x;\vartheta,r))
  \sigma(z_\beta(x';\vartheta,r))
  \rho(\dd \vartheta),
  \label{eq:levelb-shared-kernel}
\end{align}
and
\begin{equation}
  K_{\alpha\alpha}^{\mathrm{sh,head}}(x,x';\rho,r)
  :=
  \langle g_\alpha^{\mathrm{sh}}(x;\rho,r),g_\alpha^{\mathrm{sh}}(x';\rho,r)\rangle_{\R^d}.
  \label{eq:levelb-shared-head-kernel}
\end{equation}
\end{proposition}

\begin{proof}
The proof is the same differentiation argument as in \cref{prop:levelb-dynamics}.
Indeed,
\[
  \partial_c\Psi_\alpha^{\mathrm{sh}}=\sigma(z_\alpha),
  \qquad
  \partial_b\Psi_\alpha^{\mathrm{sh}}=c\,\sigma'(z_\alpha),
  \qquad
  \nabla_w\Psi_\alpha^{\mathrm{sh}}=c\,\sigma'(z_\alpha)(r_\alpha\odot x),
\]
so the neuron contribution produces \eqref{eq:levelb-shared-kernel}, while
\[
  \nabla_{r_\alpha}f_\alpha^{\mathrm{sh}}(x;\rho,r)=g_\alpha^{\mathrm{sh}}(x;\rho,r)
\]
and the head ODE gives \eqref{eq:levelb-shared-head-kernel}.
\end{proof}

\begin{corollary}[Shared output strengthens direct collectivity on the symmetric manifold]
\label{cor:levelb-shared-collapsed}
Assume that $(\rho_t,r_t)_{t\in[0,T]}$ is a regular characteristic solution of the shared-output Level-B system such that
\[
  r_{1,t}=\cdots=r_{M,t}=:r_t^\star
\]
for every $t$.
Then the outputs automatically collapse to a common predictor
\[
  f_1^{\mathrm{sh}}(x;\rho_t,r_t)=\cdots=f_M^{\mathrm{sh}}(x;\rho_t,r_t)=:f^{\mathrm{sh}}(x,t).
\]
Define
\[
  H^{\mathrm{sh}}
  :=
  \int_{\Theta^{\mathrm{sh}}}
  c^2\sigma'(z(x))\sigma'(z(x'))
  \bigl[1+(r_t^\star\odot x)^\top(r_t^\star\odot x')\bigr]
  \rho_t(\dd \vartheta),
\]
with
\[
  z(x):=w^\top(r_t^\star\odot x)+b,
\]
and
\[
  S^{\mathrm{sh}}(x,x';\rho_t,r_t^\star)
  :=
  \int_{\Theta^{\mathrm{sh}}}
  \sigma(z(x))\sigma(z(x'))
  \rho_t(\dd \vartheta).
\]
Then
\begin{equation}
  \partial_t f^{\mathrm{sh}}(x,t)
  =
  -\E_{(x',y)\sim\mathsf P}
  \Bigl[
    K_{\mathrm{eff}}^{\mathrm{sh}}(x,x';\rho_t,r_t^\star)
    \partial_1\ell(f^{\mathrm{sh}}(x',t),y)
  \Bigr],
  \label{eq:levelb-shared-keff-dynamics}
\end{equation}
where
\begin{equation}
  K_{\mathrm{eff}}^{\mathrm{sh}}
  :=
  H^{\mathrm{sh}}
  +
  \frac{K^{\mathrm{sh,head}}}{M}
  +
  S^{\mathrm{sh}},
  \label{eq:levelb-shared-keff}
\end{equation}
and
\[
  K^{\mathrm{sh,head}}
  :=
  K_{\alpha\alpha}^{\mathrm{sh,head}},
\]
which is independent of $\alpha$ by symmetry.
In particular, the direct last-layer term $S^{\mathrm{sh}}$ survives at order one on the symmetric manifold, whereas in the TabM-like Level-B formula \cref{eq:levelb-keff} the corresponding diagonal-only output term is averaged by the factor $1/M$.
\end{corollary}

\begin{proof}
Under $r_{1,t}=\cdots=r_{M,t}$, the last term in \eqref{eq:levelb-shared-kernel} is the same for every pair $(\alpha,\beta)$, namely $S^{\mathrm{sh}}$.
Therefore its row average is again $S^{\mathrm{sh}}$, rather than $S^{\mathrm{sh}}/M$.
Under the same symmetry, the hidden-feature term in \eqref{eq:levelb-shared-kernel} is also independent of the pair $(\alpha,\beta)$, so its row average is simply $H^{\mathrm{sh}}$.
The head-kernel term remains diagonal in $\alpha$ and therefore contributes $K^{\mathrm{sh,head}}/M$.
This yields \eqref{eq:levelb-shared-keff-dynamics}--\eqref{eq:levelb-shared-keff}.
The final comparison with the TabM-like Level-B model follows because the output-only contribution in \cref{eq:levelb-neuron-kernel} is diagonal in $\alpha,\beta$, so after row averaging it enters with coefficient $1/M$.
\end{proof}

\begin{remark}
\Cref{cor:levelb-shared-collapsed} gives a clean practitioner-facing comparison:
within the same semifaithful two-layer mean-field framework, independent output heads suppress the most direct last-layer route to collectivity by a factor of $M$ on the symmetric manifold.
This does not eliminate collective behavior coming from shared hidden features or global head variables,
but it does remove the strongest immediate output-level coupling mechanism.
\end{remark}

\subsection{Companion note on reduced surrogates}

To keep the present note focused on the abstract theory and the semifaithful Level-B model, the more aggressive reduced-surrogate calculations have been moved to the companion file
\[
  \texttt{tabm\_meanfield\_levelc\_surrogates.tex}.
\]
That companion note records the archived Level-C material:
compact-domain transfer of the abstract theory to a particle-friendly scalar surrogate,
explicit reduced-surrogate kernel formulas,
shared-output versus independent-output comparisons,
and one-feature small-time coupling-generation theorems.

\section{Deviation Dynamics and Practitioner Implications}

\subsection{Exact average/deviation decomposition}

The collapse result from \cref{cor:collapse} can be sharpened into an exact decomposition that is useful even away from the symmetric manifold.

\begin{proposition}[Average and deviation dynamics]
\label{prop:deviation}
Let
\[
  \bar f(x,t)=\frac1M \sum_{\alpha=1}^M f_\alpha(x;\rho_t),
  \qquad
  \delta_\alpha(x,t)=f_\alpha(x;\rho_t)-\bar f(x,t).
\]
Define
\[
  g_\beta(x',y,t)=\partial_1 \ell(f_\beta(x';\rho_t),y),
\]
the row-average kernel
\[
  \bar K_\beta(x,x';t)=\frac1M \sum_{\gamma=1}^M K_{\gamma\beta}(x,x';\rho_t),
\]
and the centered kernel
\[
  \widetilde K_{\alpha\beta}(x,x';t)=K_{\alpha\beta}(x,x';\rho_t)-\bar K_\beta(x,x';t).
\]
Then
\begin{equation}
  \partial_t \bar f(x,t)
  =
  -
  \E_{(x',y)\sim \mathsf{P}}
  \Bigl[
    \frac1M \sum_{\beta=1}^M
    \bar K_\beta(x,x';t)\, g_\beta(x',y,t)
  \Bigr],
  \label{eq:avgdyn}
\end{equation}
and, for each $\alpha$,
\begin{equation}
  \partial_t \delta_\alpha(x,t)
  =
  -
  \E_{(x',y)\sim \mathsf{P}}
  \Bigl[
    \frac1M \sum_{\beta=1}^M
    \widetilde K_{\alpha\beta}(x,x';t)\, g_\beta(x',y,t)
  \Bigr].
  \label{eq:devdyn}
\end{equation}
\end{proposition}

\begin{proof}
Average \cref{eq:output-dynamics} over $\alpha$:
\begin{align*}
  \partial_t \bar f(x,t)
  &=
  \frac1M \sum_{\alpha=1}^M \partial_t f_\alpha(x;\rho_t)
  \\
  &=
  -
  \E_{(x',y)}
  \Bigl[
    \frac1M \sum_{\beta=1}^M
    \Bigl(
      \frac1M \sum_{\alpha=1}^M K_{\alpha\beta}(x,x';\rho_t)
    \Bigr)
    g_\beta(x',y,t)
  \Bigr],
\end{align*}
which is \eqref{eq:avgdyn}.

Now subtract \eqref{eq:avgdyn} from \cref{eq:output-dynamics}:
\begin{align*}
  \partial_t \delta_\alpha(x,t)
  &=
  -
  \E_{(x',y)}
  \Bigl[
    \frac1M \sum_{\beta=1}^M
    \bigl(
      K_{\alpha\beta}(x,x';\rho_t)-\bar K_\beta(x,x';t)
    \bigr)
    g_\beta(x',y,t)
  \Bigr],
\end{align*}
which is \eqref{eq:devdyn}.
\end{proof}

\begin{corollary}[Only centered row-variation can drive headwise differences]
\label{cor:rowvariation}
The deterministic deviation dynamics depends only on the centered kernel $\widetilde K$.
In particular, if
\[
  K_{\alpha\beta}(x,x';\rho_t)
  \quad
  \text{is independent of }\alpha
  \quad
  \text{for every }\beta,x,x',
\]
then
\[
  \partial_t \delta_\alpha(x,t)=0
  \qquad
  \forall \alpha,x.
\]
\end{corollary}

\begin{proof}
If $K_{\alpha\beta}$ is independent of $\alpha$, then
\[
  K_{\alpha\beta}(x,x';\rho_t)=\bar K_\beta(x,x';t)
  \qquad \forall \alpha,\beta,x,x',
\]
hence $\widetilde K_{\alpha\beta}\equiv 0$ and \eqref{eq:devdyn} gives the claim.
\end{proof}

\begin{remark}
\Cref{prop:deviation,cor:rowvariation} give a sharp practitioner-facing summary:
in deterministic mean-field, headwise diversity is not driven by the common part of the kernel,
but only by the row-to-row asymmetry of the ensemble kernel.
Exact symmetry removes that asymmetry completely.
\end{remark}

\subsection{A local transverse criterion from the frozen-kernel tangent model}

The nonlinear stability theorem above controls the full PDE in Wasserstein distance.
To obtain a sharper spectral criterion, it is useful to freeze the kernel along a symmetric reference trajectory
and linearize only the loss term.
This is not the full feature-learning linearization of the mean-field PDE,
but it cleanly identifies the operator that governs disagreement directions.

\begin{proposition}[Average and disagreement modes decouple in the frozen-kernel tangent model]
\label{prop:frozen-tangent}
Assume squared loss
\[
  \ell(z,y)=\frac12 (z-y)^2.
\]
Let $\rho_t^\star$ be a permutation-invariant reference solution of \eqref{eq:pde},
and let
\[
  f^\star(x,t)=f_1(x;\rho_t^\star).
\]
Write
\[
  K_{\alpha\alpha}(x,x';\rho_t^\star)=K_{\mathrm d}(x,x';t),
  \qquad
  K_{\alpha\beta}(x,x';\rho_t^\star)=K_{\mathrm o}(x,x';t)
  \quad (\alpha\neq\beta).
\]
Let $\mathsf P_X$ denote the marginal of $\mathsf P$ on $\mathcal X$.
Consider the frozen-kernel tangent system
\begin{equation}
  \partial_t \varepsilon_\alpha(x,t)
  =
  -
  \int_{\mathcal X}
  \frac1M \sum_{\beta=1}^M
  \mathcal K_{\alpha\beta}(x,x';t)\,
  \varepsilon_\beta(x',t)\,
  \mathsf P_X(\dd x'),
  \label{eq:frozen-tangent}
\end{equation}
where
\[
  \mathcal K_{\alpha\alpha}(x,x';t)=K_{\mathrm d}(x,x';t),
  \qquad
  \mathcal K_{\alpha\beta}(x,x';t)=K_{\mathrm o}(x,x';t)
  \quad (\alpha\neq\beta).
\]
Define
\[
  \bar \varepsilon(x,t)=\frac1M \sum_{\alpha=1}^M \varepsilon_\alpha(x,t),
  \qquad
  \Delta_\alpha(x,t)=\varepsilon_\alpha(x,t)-\bar \varepsilon(x,t).
\]
Then the average mode and the disagreement modes decouple:
\begin{equation}
  \partial_t \bar \varepsilon(x,t)
  =
  -
  \int_{\mathcal X}
  K_{\mathrm{eff}}(x,x';t)\,
  \bar \varepsilon(x',t)\,
  \mathsf P_X(\dd x'),
  \label{eq:avg-tangent}
\end{equation}
with
\[
  K_{\mathrm{eff}}=\frac{K_{\mathrm d}+(M-1)K_{\mathrm o}}{M},
\]
and, for every $\alpha$,
\begin{equation}
  \partial_t \Delta_\alpha(x,t)
  =
  -
  \frac1M
  \int_{\mathcal X}
  \bigl(
    K_{\mathrm d}(x,x';t)-K_{\mathrm o}(x,x';t)
  \bigr)\,
  \Delta_\alpha(x',t)\,
  \mathsf P_X(\dd x').
  \label{eq:disc-tangent}
\end{equation}
\end{proposition}

\begin{proof}
Average \eqref{eq:frozen-tangent} over $\alpha$.
Because
\[
  \frac1M \sum_{\alpha=1}^M \mathcal K_{\alpha\beta}(x,x';t)
  =
  \frac{K_{\mathrm d}(x,x';t)+(M-1)K_{\mathrm o}(x,x';t)}{M}
  =
  K_{\mathrm{eff}}(x,x';t),
\]
we obtain \eqref{eq:avg-tangent}.

Now subtract \eqref{eq:avg-tangent} from \eqref{eq:frozen-tangent}.
Using
\[
  \varepsilon_\beta=\bar \varepsilon+\Delta_\beta,
  \qquad
  \sum_{\beta=1}^M \Delta_\beta=0,
\]
the common part proportional to $\bar \varepsilon$ cancels.
For the disagreement part, the off-diagonal sum reduces to
\[
  \sum_{\beta\neq\alpha} K_{\mathrm o}(x,x';t)\Delta_\beta(x',t)
  =
  -K_{\mathrm o}(x,x';t)\Delta_\alpha(x',t),
\]
because the disagreement components sum to zero.
Therefore
\[
  \frac1M \sum_{\beta=1}^M
  \mathcal K_{\alpha\beta}(x,x';t)\Delta_\beta(x',t)
  =
  \frac{K_{\mathrm d}(x,x';t)-K_{\mathrm o}(x,x';t)}{M}
  \Delta_\alpha(x',t),
\]
which is exactly \eqref{eq:disc-tangent}.
\end{proof}

\begin{corollary}[A sufficient frozen-kernel contraction criterion]
\label{cor:frozen-contraction}
Assume the hypotheses of \cref{prop:frozen-tangent}.
Suppose that for every $t\in[0,T]$ there exists $\kappa_t\ge 0$ such that
\begin{equation}
  \int_{\mathcal X\times \mathcal X}
  \varphi(x)
  \bigl(
    K_{\mathrm d}(x,x';t)-K_{\mathrm o}(x,x';t)
  \bigr)
  \varphi(x')
  \,\mathsf P_X(\dd x)\mathsf P_X(\dd x')
  \ge
  \kappa_t \|\varphi\|_{L^2(\mathsf P_X)}^2
  \label{eq:kappa-cond}
\end{equation}
for every $\varphi \in L^2(\mathsf P_X)$.
Then the disagreement energy satisfies
\begin{equation}
  \sum_{\alpha=1}^M
  \|\Delta_\alpha(\cdot,t)\|_{L^2(\mathsf P_X)}^2
  \le
  \exp\Bigl(
    -\frac{2}{M}\int_0^t \kappa_s\,\dd s
  \Bigr)
  \sum_{\alpha=1}^M
  \|\Delta_\alpha(\cdot,0)\|_{L^2(\mathsf P_X)}^2.
  \label{eq:frozen-decay}
\end{equation}
\end{corollary}

\begin{proof}
By \cref{prop:frozen-tangent}, each disagreement mode solves \eqref{eq:disc-tangent}.
Hence
\begin{align*}
  \frac12 \frac{\dd}{\dd t}
  \|\Delta_\alpha(\cdot,t)\|_{L^2(\mathsf P_X)}^2
  &=
  -\frac1M
  \int_{\mathcal X\times \mathcal X}
  \Delta_\alpha(x,t)
  \bigl(
    K_{\mathrm d}(x,x';t)-K_{\mathrm o}(x,x';t)
  \bigr)
  \Delta_\alpha(x',t)
  \,\mathsf P_X(\dd x)\mathsf P_X(\dd x')
  \\
  &\le
  -\frac{\kappa_t}{M}
  \|\Delta_\alpha(\cdot,t)\|_{L^2(\mathsf P_X)}^2
\end{align*}
by \eqref{eq:kappa-cond}.
Sum over $\alpha$ and apply Gr\"onwall's inequality to obtain \eqref{eq:frozen-decay}.
\end{proof}

\begin{remark}
\Cref{prop:frozen-tangent,cor:frozen-contraction} isolate a sharp local message:
in the kernel-linearized transverse dynamics,
the quantity that matters is neither $K_{\mathrm d}$ nor $K_{\mathrm o}$ separately,
but the \emph{difference} $K_{\mathrm d}-K_{\mathrm o}$.
Large positive $K_{\mathrm d}-K_{\mathrm o}$ contracts disagreement;
if this operator becomes weak or indefinite,
then finite-width asymmetries can persist or be amplified.
\end{remark}

\subsection{Level-B transverse operators: common-mode and disagreement-mode effects differ}

The preceding shared-output comparison concerns the \emph{common predictor} on the symmetric manifold.
For practitioner-facing interpretation, it is equally important to understand the local disagreement operator.
The next result shows that the independent-head and shared-output architectures behave differently here as well,
but in a subtler way:
independent output heads weaken average-mode collectivity while simultaneously adding extra diagonal self-coupling terms in the frozen transverse dynamics.

\begin{proposition}[Frozen transverse kernels for Level B and its shared-output proxy]
\label{prop:levelb-transverse}
Assume squared loss and let $(\rho_t^\star,r_t^\star)_{t\in[0,T]}$ be a regular symmetric characteristic solution of the untruncated TabM-like Level-B system, in the sense that
\[
  r_{1,t}^\star=\cdots=r_{M,t}^\star=:r_t^\star
\]
for every $t$ and $\rho_t^\star$ is exchangeable across $(a_\alpha)_{\alpha=1}^M$.
Define
\[
  z_t(x;\vartheta):=w^\top(r_t^\star\odot x)+b,
\]
\[
  q_t(x,x';\vartheta)
  :=
  \sigma'(z_t(x;\vartheta))
  \sigma'(z_t(x';\vartheta))
  \bigl[
    1+(r_t^\star\odot x)^\top(r_t^\star\odot x')
  \bigr],
\]
and the kernels
\begin{align}
  S_t^{\mathrm{out}}(x,x')
  &:=
  \int_\Theta
  \sigma(z_t(x;\vartheta))
  \sigma(z_t(x';\vartheta))
  \,\rho_t^\star(\dd \vartheta),
  \label{eq:levelb-transverse-sout}
  \\
  M_t^{\mathrm{gap}}(x,x')
  &:=
  \frac12
  \int_\Theta
  (a_1-a_2)^2
  q_t(x,x';\vartheta)
  \,\rho_t^\star(\dd \vartheta),
  \label{eq:levelb-transverse-mgap}
  \\
  H_t^{\mathrm{ind}}(x,x')
  &:=
  K^{\mathrm{head}}(x,x';\rho_t^\star,r_t^\star).
  \label{eq:levelb-transverse-hind}
\end{align}
Then the frozen-kernel disagreement system from \cref{prop:frozen-tangent} is
\begin{equation}
  \partial_t \Delta_\alpha(x,t)
  =
  -
  \frac1M
  \int_{\mathcal X}
  T_t^{\mathrm{ind}}(x,x')\,
  \Delta_\alpha(x',t)\,
  \mathsf P_X(\dd x'),
  \label{eq:levelb-transverse-ind}
\end{equation}
with
\begin{equation}
  T_t^{\mathrm{ind}}
  =
  S_t^{\mathrm{out}}
  +
  M_t^{\mathrm{gap}}
  +
  H_t^{\mathrm{ind}}.
  \label{eq:levelb-transverse-ind-kernel}
\end{equation}

Now let $(\rho_t^{\mathrm{sh},\star},r_t^{\mathrm{sh},\star})_{t\in[0,T]}$ be a regular symmetric characteristic solution of the shared-output Level-B proxy, with
\[
  r_{1,t}^{\mathrm{sh},\star}=\cdots=r_{M,t}^{\mathrm{sh},\star}=:r_t^{\mathrm{sh},\star}.
\]
Define
\[
  H_t^{\mathrm{sh}}(x,x')
  :=
  K^{\mathrm{sh,head}}(x,x';\rho_t^{\mathrm{sh},\star},r_t^{\mathrm{sh},\star}).
\]
Then its frozen-kernel disagreement system is
\begin{equation}
  \partial_t \Delta_\alpha^{\mathrm{sh}}(x,t)
  =
  -
  \frac1M
  \int_{\mathcal X}
  T_t^{\mathrm{sh}}(x,x')\,
  \Delta_\alpha^{\mathrm{sh}}(x',t)\,
  \mathsf P_X(\dd x'),
  \label{eq:levelb-transverse-sh}
\end{equation}
with
\begin{equation}
  T_t^{\mathrm{sh}}=H_t^{\mathrm{sh}}.
  \label{eq:levelb-transverse-sh-kernel}
\end{equation}

Moreover, each of the kernels
\[
  S_t^{\mathrm{out}},
  \qquad
  M_t^{\mathrm{gap}},
  \qquad
  H_t^{\mathrm{ind}},
  \qquad
  H_t^{\mathrm{sh}}
\]
is positive semidefinite on $L^2(\mathsf P_X)$.
\end{proposition}

\begin{proof}
For the TabM-like Level-B model, \cref{prop:frozen-tangent} applies to the total kernel
\[
  \mathcal K_{\alpha\beta}^{\mathrm{ind}}
  :=
  K_{\alpha\beta}^{\mathrm{neu}}
  +
  \delta_{\alpha\beta}K^{\mathrm{head}}_{\alpha\alpha}.
\]
Hence the disagreement operator is governed by
\[
  K_{\mathrm d}^{\mathrm{tot}}-K_{\mathrm o}^{\mathrm{tot}}
  =
  \bigl(K_{\mathrm d}^{\mathrm{neu}}-K_{\mathrm o}^{\mathrm{neu}}\bigr)
  +
  K^{\mathrm{head}},
\]
where, on the symmetric reference path, $K^{\mathrm{head}}=H_t^{\mathrm{ind}}$ is independent of the head index.
By \cref{eq:levelb-neuron-kernel},
\[
  K_{\mathrm d}^{\mathrm{neu}}-K_{\mathrm o}^{\mathrm{neu}}
  =
  S_t^{\mathrm{out}}
  +
  \int_\Theta
  (a_1^2-a_1a_2)\,q_t(x,x';\vartheta)\,\rho_t^\star(\dd \vartheta).
\]
Exchangeability implies
\[
  \int_\Theta a_1^2 q_t\,\rho_t^\star(\dd \vartheta)
  =
  \int_\Theta a_2^2 q_t\,\rho_t^\star(\dd \vartheta),
\]
so
\[
  \int_\Theta
  (a_1^2-a_1a_2)\,q_t\,\rho_t^\star(\dd \vartheta)
  =
  \frac12
  \int_\Theta
  (a_1-a_2)^2 q_t\,\rho_t^\star(\dd \vartheta)
  =
  M_t^{\mathrm{gap}}.
\]
This proves \eqref{eq:levelb-transverse-ind}--\eqref{eq:levelb-transverse-ind-kernel}.

For the shared-output proxy, the total frozen kernel is
\[
  \mathcal K_{\alpha\beta}^{\mathrm{sh}}
  :=
  K_{\alpha\beta}^{\mathrm{sh}}
  +
  \delta_{\alpha\beta}K_{\alpha\alpha}^{\mathrm{sh,head}}.
\]
On the symmetric reference path, \cref{eq:levelb-shared-kernel} is independent of the pair $(\alpha,\beta)$.
Therefore its contribution cancels in
\[
  K_{\mathrm d}^{\mathrm{sh,tot}}-K_{\mathrm o}^{\mathrm{sh,tot}},
\]
and only the diagonal head-kernel term remains.
This gives \eqref{eq:levelb-transverse-sh}--\eqref{eq:levelb-transverse-sh-kernel}.

It remains to prove positive semidefiniteness.
The kernel $S_t^{\mathrm{out}}$ is of the standard feature-product form
\[
  S_t^{\mathrm{out}}(x,x')
  =
  \int_\Theta
  \phi_t(x;\vartheta)\phi_t(x';\vartheta)\,\rho_t^\star(\dd \vartheta),
  \qquad
  \phi_t(x;\vartheta):=\sigma(z_t(x;\vartheta)),
\]
so it is positive semidefinite.
Likewise $H_t^{\mathrm{ind}}$ and $H_t^{\mathrm{sh}}$ are inner-product kernels by definition.

For $M_t^{\mathrm{gap}}$, define
\[
  \Phi_t(x;\vartheta)
  :=
  \Bigl(
    \sigma'(z_t(x;\vartheta)),
    \ \sigma'(z_t(x;\vartheta))(r_t^\star\odot x)
  \Bigr)\in \R^{1+d}.
\]
Then
\[
  q_t(x,x';\vartheta)=\langle \Phi_t(x;\vartheta),\Phi_t(x';\vartheta)\rangle_{\R^{1+d}},
\]
and therefore
\[
  M_t^{\mathrm{gap}}(x,x')
  =
  \frac12
  \int_\Theta
  (a_1-a_2)^2
  \langle \Phi_t(x;\vartheta),\Phi_t(x';\vartheta)\rangle
  \,\rho_t^\star(\dd \vartheta),
\]
which is again positive semidefinite because the weight $(a_1-a_2)^2/2$ is nonnegative.
\end{proof}

\begin{remark}
\label{rem:levelb-two-sided}
\Cref{cor:levelb-shared-collapsed,prop:levelb-transverse} reveal a two-sided role for independent output heads.
On the symmetric manifold, they weaken the most direct last-layer contribution to the \emph{common predictor}, because the output-only term is reduced by a factor of $1/M$ in \cref{eq:levelb-keff}.
But in the local frozen disagreement dynamics they add the positive-semidefinite kernels $S_t^{\mathrm{out}}$ and $M_t^{\mathrm{gap}}$ that are absent from the shared-output proxy.
So independent heads do \emph{not} by themselves produce an NTK-like independent regime inside deterministic mean-field.
If substantial head diversity is observed in practice, then it must come from finite-width fluctuations, explicit symmetry breaking, or genuinely nonlocal feature-learning effects rather than from deterministic leading-order transverse dynamics alone.
\end{remark}

\begin{corollary}[Mode allocation of the output-only last-layer term]
\label{cor:levelb-mode-allocation}
Under the hypotheses of \cref{cor:levelb-collapsed,cor:levelb-shared-collapsed,prop:levelb-transverse},
the corresponding output-only feature-product term
\[
  S(x,x')
  :=
  \int
  \sigma(z(x))\sigma(z(x'))
  \,\dd \nu
\]
where $\nu$ denotes the relevant symmetric law in the model under consideration,
appears in different places depending on the output-head design:
\begin{enumerate}[leftmargin=2em]
  \item in the TabM-like independent-head model, it contributes to the common-mode dynamics with coefficient $1/M$ through \cref{eq:levelb-keff}, and to the frozen disagreement dynamics with coefficient $1/M$ through \cref{eq:levelb-transverse-ind};
  \item in the shared-output proxy, it contributes to the common-mode dynamics with coefficient $1$ through \cref{eq:levelb-shared-keff}, and it cancels completely from the frozen disagreement dynamics.
\end{enumerate}
So independent heads do not simply ``add diversity'' or ``remove collectivity'':
they reallocate the most direct output-level term away from the common predictor and into the local disagreement operator.
\end{corollary}

\begin{proof}
The statement is immediate from \cref{eq:levelb-keff,eq:levelb-transverse-ind-kernel,eq:levelb-shared-keff,eq:levelb-transverse-sh-kernel}.
For the independent-head model, the output-only term is diagonal in the ensemble index, so row averaging on the symmetric manifold produces the factor $1/M$, while the disagreement dynamics inherits the same term through $K_{\mathrm d}^{\mathrm{tot}}-K_{\mathrm o}^{\mathrm{tot}}$ and the prefactor $1/M$ in \cref{eq:levelb-transverse-ind}.
For the shared-output proxy, the corresponding feature-product term is identical for every pair $(\alpha,\beta)$ on the symmetric manifold, so it survives unchanged in the row average but cancels in the diagonal-minus-off-diagonal difference.
\end{proof}

\subsection{A fluctuation template from the frozen tangent system}

The full finite-width central-limit theory for Level B remains open.
Nevertheless, the frozen tangent system already gives a mathematically clean template for how
\[
  m^{-1/2}\text{-scale fluctuations}
\]
should split into common-mode and disagreement-mode components.
This is the most direct bridge between the deterministic mean-field picture and the practically observed residual diversity at large but finite width.

\begin{proposition}[Covariance dynamics for frozen average and disagreement modes]
\label{prop:frozen-covariance}
Assume the hypotheses of \cref{prop:frozen-tangent}.
Let
\[
  (\varepsilon_\alpha(\cdot,t))_{\alpha=1}^M
\]
be a square-integrable random solution of the frozen tangent system \eqref{eq:frozen-tangent}.
Define the random average and disagreement fields
\[
  \bar \varepsilon(\cdot,t)=\frac1M\sum_{\alpha=1}^M \varepsilon_\alpha(\cdot,t),
  \qquad
  \Delta_\alpha(\cdot,t)=\varepsilon_\alpha(\cdot,t)-\bar\varepsilon(\cdot,t).
\]
Let $A_t^{\mathrm{avg}}$ and $A_t^{\mathrm{tr}}$ be the integral operators on $L^2(\mathsf P_X)$
given by
\[
  (A_t^{\mathrm{avg}}\varphi)(x)
  :=
  \int_{\mathcal X}
  K_{\mathrm{eff}}(x,x';t)\,\varphi(x')\,\mathsf P_X(\dd x'),
\]
\[
  (A_t^{\mathrm{tr}}\varphi)(x)
  :=
  \frac1M\int_{\mathcal X}
  \bigl(K_{\mathrm d}(x,x';t)-K_{\mathrm o}(x,x';t)\bigr)\,
  \varphi(x')\,\mathsf P_X(\dd x').
\]
For $\phi,\psi\in L^2(\mathsf P_X)$, define the covariance bilinear forms
\[
  \langle \phi, C_t^{\mathrm{avg}}\psi\rangle
  :=
  \E\bigl[\langle \phi,\bar\varepsilon_t\rangle \langle \psi,\bar\varepsilon_t\rangle\bigr],
  \qquad
  \langle \phi, C_{\alpha,t}^{\mathrm{tr}}\psi\rangle
  :=
  \E\bigl[\langle \phi,\Delta_{\alpha,t}\rangle \langle \psi,\Delta_{\alpha,t}\rangle\bigr].
\]
Then
\begin{equation}
  \partial_t C_t^{\mathrm{avg}}
  =
  -A_t^{\mathrm{avg}} C_t^{\mathrm{avg}}
  -
  C_t^{\mathrm{avg}} (A_t^{\mathrm{avg}})^*,
  \label{eq:frozen-cov-avg}
\end{equation}
and, for every $\alpha$,
\begin{equation}
  \partial_t C_{\alpha,t}^{\mathrm{tr}}
  =
  -A_t^{\mathrm{tr}} C_{\alpha,t}^{\mathrm{tr}}
  -
  C_{\alpha,t}^{\mathrm{tr}} (A_t^{\mathrm{tr}})^*.
  \label{eq:frozen-cov-tr}
\end{equation}
If, in addition, the initial law is exchangeable across heads, then
\[
  C_{1,t}^{\mathrm{tr}}=\cdots=C_{M,t}^{\mathrm{tr}}
\]
for all $t$.
\end{proposition}

\begin{proof}
By \cref{prop:frozen-tangent}, the average field solves
\[
  \partial_t \bar\varepsilon_t = -A_t^{\mathrm{avg}}\bar\varepsilon_t,
\]
while each disagreement field solves
\[
  \partial_t \Delta_{\alpha,t} = -A_t^{\mathrm{tr}}\Delta_{\alpha,t}.
\]
Fix $\phi,\psi\in L^2(\mathsf P_X)$.
Then
\[
  \frac{\dd}{\dd t}\langle \phi,\bar\varepsilon_t\rangle
  =
  -\langle (A_t^{\mathrm{avg}})^*\phi,\bar\varepsilon_t\rangle.
\]
Therefore
\begin{align*}
  \frac{\dd}{\dd t}
  \E\bigl[\langle \phi,\bar\varepsilon_t\rangle \langle \psi,\bar\varepsilon_t\rangle\bigr]
  &=
  -\E\bigl[\langle (A_t^{\mathrm{avg}})^*\phi,\bar\varepsilon_t\rangle \langle \psi,\bar\varepsilon_t\rangle\bigr]
  \\
  &\quad
  -
  \E\bigl[\langle \phi,\bar\varepsilon_t\rangle \langle (A_t^{\mathrm{avg}})^*\psi,\bar\varepsilon_t\rangle\bigr],
\end{align*}
which is exactly \eqref{eq:frozen-cov-avg}.
The same computation with $\Delta_{\alpha,t}$ proves \eqref{eq:frozen-cov-tr}.

If the initial law is exchangeable across heads, then so is the frozen tangent system itself.
Hence the laws of the disagreement fields remain exchangeable, which implies equality of the covariance forms $C_{\alpha,t}^{\mathrm{tr}}$ across $\alpha$.
\end{proof}

\begin{corollary}[Frozen fluctuation decay under transverse coercivity]
\label{cor:frozen-fluctuation-decay}
Under the hypotheses of \cref{cor:frozen-contraction,prop:frozen-covariance},
assume that the disagreement fluctuations are centered and exchangeable across heads.
Then, for every $\alpha$,
\begin{equation}
  \mathrm{Tr}\, C_{\alpha,t}^{\mathrm{tr}}
  \le
  \exp\Bigl(
    -\frac{2}{M}\int_0^t \kappa_s\,\dd s
  \Bigr)
  \mathrm{Tr}\, C_{\alpha,0}^{\mathrm{tr}},
  \label{eq:frozen-fluctuation-trace}
\end{equation}
whenever the covariance operators are trace class.
\end{corollary}

\begin{proof}
Apply \cref{cor:frozen-contraction} pathwise to each disagreement realization.
Taking expectations gives
\[
  \E\|\Delta_{\alpha,t}\|_{L^2(\mathsf P_X)}^2
  \le
  \exp\Bigl(
    -\frac{2}{M}\int_0^t \kappa_s\,\dd s
  \Bigr)
  \E\|\Delta_{\alpha,0}\|_{L^2(\mathsf P_X)}^2.
\]
For centered square-integrable random fields, the left-hand side equals $\mathrm{Tr}\, C_{\alpha,t}^{\mathrm{tr}}$ and the right-hand side equals $\mathrm{Tr}\, C_{\alpha,0}^{\mathrm{tr}}$.
\end{proof}

\begin{proposition}[Pairwise disagreement is exactly the transverse fluctuation energy]
\label{prop:frozen-pairwise}
Under the hypotheses of \cref{prop:frozen-covariance}, assume in addition that
\[
  \sum_{\alpha=1}^M \Delta_{\alpha,t}(x)=0
  \qquad
  \text{for every }x,t,
\]
and that the law of $(\Delta_{\alpha,t})_{\alpha=1}^M$ is exchangeable across heads.
Define the frozen pairwise disagreement energy
\[
  \mathcal D_t
  :=
  \frac{1}{M(M-1)}
  \sum_{\alpha\neq\beta}
  \E\bigl[\|\Delta_{\alpha,t}-\Delta_{\beta,t}\|_{L^2(\mathsf P_X)}^2\bigr].
\]
Then
\begin{equation}
  \mathcal D_t
  =
  \frac{2M}{M-1}\,
  \mathrm{Tr}\, C_{1,t}^{\mathrm{tr}}.
  \label{eq:frozen-pairwise-trace}
\end{equation}
\end{proposition}

\begin{proof}
Because \(\sum_{\alpha=1}^M \Delta_{\alpha,t}=0\), the standard identity
\[
  \sum_{\alpha<\beta}
  \|\Delta_{\alpha,t}-\Delta_{\beta,t}\|_{L^2(\mathsf P_X)}^2
  =
  M\sum_{\alpha=1}^M
  \|\Delta_{\alpha,t}\|_{L^2(\mathsf P_X)}^2
\]
holds pathwise.
Doubling both sides yields
\[
  \sum_{\alpha\neq\beta}
  \|\Delta_{\alpha,t}-\Delta_{\beta,t}\|_{L^2(\mathsf P_X)}^2
  =
  2M\sum_{\alpha=1}^M
  \|\Delta_{\alpha,t}\|_{L^2(\mathsf P_X)}^2.
\]
Divide by \(M(M-1)\) and take expectations:
\[
  \mathcal D_t
  =
  \frac{2}{M-1}
  \sum_{\alpha=1}^M
  \E\bigl[\|\Delta_{\alpha,t}\|_{L^2(\mathsf P_X)}^2\bigr].
\]
Exchangeability gives
\[
  \E\bigl[\|\Delta_{\alpha,t}\|_{L^2(\mathsf P_X)}^2\bigr]
  =
  \E\bigl[\|\Delta_{1,t}\|_{L^2(\mathsf P_X)}^2\bigr]
  =
  \mathrm{Tr}\, C_{1,t}^{\mathrm{tr}},
\]
which proves \eqref{eq:frozen-pairwise-trace}.
\end{proof}

\begin{corollary}[Frozen pairwise disagreement decays with the transverse covariance]
\label{cor:frozen-pairwise-decay}
Under the hypotheses of \cref{cor:frozen-fluctuation-decay,prop:frozen-pairwise},
\begin{equation}
  \mathcal D_t
  \le
  \exp\Bigl(
    -\frac{2}{M}\int_0^t \kappa_s\,\dd s
  \Bigr)
  \mathcal D_0.
  \label{eq:frozen-pairwise-decay}
\end{equation}
\end{corollary}

\begin{proof}
Combine \cref{eq:frozen-pairwise-trace} with \cref{eq:frozen-fluctuation-trace}.
\end{proof}

\begin{corollary}[Exact $m^{-1}$ width law in the frozen CLT model]
\label{cor:frozen-width-law}
Under the hypotheses of \cref{prop:frozen-pairwise}, fix a deterministic common predictor $f(\cdot,t)$ and,
for each width $m\ge 1$, define the auxiliary frozen-CLT outputs
\[
  \widetilde f_\alpha^{(m)}(x,t)
  :=
  f(x,t)+m^{-1/2}\Delta_\alpha(x,t).
\]
Their pairwise disagreement metric
\[
  \widetilde D_m(t)
  :=
  \frac{1}{M(M-1)}
  \sum_{\alpha\neq\beta}
  \E_x\Bigl[\bigl(\widetilde f_\alpha^{(m)}(x,t)-\widetilde f_\beta^{(m)}(x,t)\bigr)^2\Bigr]
\]
satisfies the exact identity
\begin{equation}
  \widetilde D_m(t)
  =
  \frac{1}{m}\,\mathcal D_t
  =
  \frac{2M}{M-1}\,
  \frac{1}{m}\,
  \mathrm{Tr}\,C_{1,t}^{\mathrm{tr}}.
  \label{eq:frozen-width-law}
\end{equation}
If, in addition, the transverse coercivity condition of \cref{cor:frozen-contraction} holds, then
\begin{equation}
  \widetilde D_m(t)
  \le
  \exp\Bigl(
    -\frac{2}{M}\int_0^t \kappa_s\,\dd s
  \Bigr)
  \widetilde D_m(0).
  \label{eq:frozen-width-law-decay}
\end{equation}
In particular, for every fixed time $t$, the frozen CLT model predicts an exact inverse-width law:
doubling $m$ halves the pairwise disagreement.
\end{corollary}

\begin{proof}
By definition,
\[
  \widetilde f_\alpha^{(m)}-\widetilde f_\beta^{(m)}
  =
  m^{-1/2}(\Delta_\alpha-\Delta_\beta),
\]
so
\[
  \widetilde D_m(t)
  =
  \frac{1}{m}\,
  \frac{1}{M(M-1)}
  \sum_{\alpha\neq\beta}
  \E_x\bigl[(\Delta_\alpha(x,t)-\Delta_\beta(x,t))^2\bigr]
  =
  \frac{1}{m}\,\mathcal D_t.
\]
Now apply \cref{eq:frozen-pairwise-trace} to obtain \eqref{eq:frozen-width-law}, and combine this with \cref{eq:frozen-pairwise-decay} to obtain \eqref{eq:frozen-width-law-decay}.
\end{proof}

\begin{remark}[What this means for finite-width diversity]
\label{rem:frozen-fluctuation}
\Cref{prop:frozen-covariance,cor:frozen-fluctuation-decay,prop:frozen-pairwise,cor:frozen-pairwise-decay,cor:frozen-width-law} do not yet prove a full Level-B central-limit theorem.
But they already identify the operator that any Gaussian or CLT-scale transverse fluctuation would have to obey after linearization.
So the most natural fluctuation picture is:
\begin{enumerate}[leftmargin=2em]
  \item deterministic mean-field gives the common collapsed predictor,
  \item finite-width diversity lives at the $m^{-1/2}$ scale around it,
  \item its transverse covariance is locally damped by the same operator that appears in \cref{prop:frozen-tangent},
  \item the experimentally observable pairwise disagreement is exactly proportional, in the frozen fluctuation model, to the transverse covariance trace,
  \item within the same frozen CLT template, pairwise disagreement obeys an exact $m^{-1}$ width law,
  \item hence persistent practical diversity must come either from weak transverse coercivity, from continual stochastic injection, or from genuinely nonlinear feature-learning effects that are invisible to the frozen tangent model.
\end{enumerate}
This is precisely why width sweeps and small-perturbation diagnostics are so informative experimentally.
\end{remark}

\begin{remark}[Formal CLT consequence for pairwise disagreement]
\label{rem:formal-clt-disagreement}
Suppose, purely formally, that along a symmetric Level-B trajectory one has a fluctuation expansion
\[
  f_\alpha^m(x,t)
  =
  f(x,t)+m^{-1/2}\xi_\alpha(x,t)+o(m^{-1/2})
\]
in a regime where the leading fluctuation field is described by the frozen tangent system.
Then the pairwise disagreement metric
\[
  D_m(t)
  :=
  \frac{1}{M(M-1)}\sum_{\alpha\neq\beta}
  \E_x[(f_\alpha^m(x,t)-f_\beta^m(x,t))^2]
\]
admits the formal approximation
\[
  D_m(t)
  =
  \frac{1}{m}
  \frac{1}{M(M-1)}\sum_{\alpha\neq\beta}
  \E_x[(\xi_\alpha(x,t)-\xi_\beta(x,t))^2]
  +
  o(m^{-1}).
\]
Under the transverse coercivity condition of \cref{cor:frozen-contraction},
the leading coefficient is then damped at the same exponential rate as in \cref{eq:frozen-fluctuation-trace}.
So, at the level of formal asymptotics, the theory predicts an \emph{$m^{-1}$ law for pairwise disagreement variance}, with a prefactor controlled by the transverse covariance dynamics.
\end{remark}

\begin{remark}[A practical falsification test for the frozen CLT picture]
\label{rem:frozen-loglog}
\Cref{cor:frozen-width-law} suggests a very concrete experimental diagnostic:
on a log-log plot of pairwise disagreement versus width, the frozen fluctuation template predicts slope $-1$ once the network is wide enough for the CLT-scale regime to dominate.
Therefore:
\begin{enumerate}[leftmargin=2em]
  \item a slope close to $-1$ supports the view that residual diversity is mostly a passive finite-width fluctuation around the collapsed deterministic predictor,
  \item a markedly shallower slope indicates either order-one asymmetry, continual stochastic forcing, or nonlinear amplification beyond the frozen tangent approximation,
  \item a steeper initial slope followed by saturation is consistent with a crossover between a transient contraction regime and a noise floor induced by optimization or regularization.
\end{enumerate}
This is one of the most direct practitioner-facing predictions available from the current theory.
\end{remark}

\subsection{Companion note on reduced-surrogate small-time analyses}

The one-feature and reduced-surrogate small-time coupling-generation calculations have been moved to
\[
  \texttt{tabm\_meanfield\_levelc\_surrogates.tex}.
\]
They remain useful as auxiliary mechanism calculations, but they are no longer part of the main Level-B note.

\subsection{Heuristic lessons beyond the deterministic PDE}

\begin{remark}[Why mean-field collapse and NTK diversity do not contradict each other]
\label{rem:mf-vs-ntk}
The deterministic collapse of the fixed-$M$ mean-field PDE should not be confused with the infinite-width NTK/GP limit.
The two limits keep different pieces of randomness alive.

In the deterministic mean-field scaling, the empirical measure self-averages and converges to a deterministic law.
As a result, symmetric initialization forces
\[
  f_1=\cdots=f_M
\]
at leading order.

By contrast, in the NTK/GP scaling, initialization randomness survives at order one as width tends to infinity.
This is exactly why the NTK theory of embedded ensembles can still distinguish independent and collective regimes
\cite{embedded-ensembles}.
So the correct synthesis is:
\begin{enumerate}[leftmargin=2em]
  \item deterministic feature-learning mean-field predicts collapse of the leading-order transport dynamics,
  \item NTK/GP theory can still retain order-one random diversity,
  \item practical finite networks may interpolate between these two pictures.
\end{enumerate}
\end{remark}

\begin{remark}[Finite-width origin of diversity under mean-field scaling]
\label{rem:finite-width}
If one informally writes the width-$N$ empirical measure as
\[
  \rho_t^N = \rho_t + N^{-1/2}\eta_t + \cdots,
\]
then, around a symmetric deterministic solution $\rho_t$, the corresponding outputs admit an expansion of the form
\[
  f_\alpha^N(x,t)=f(x,t)+N^{-1/2}\xi_\alpha(x,t)+\cdots.
\]
This suggests the following heuristic:
\begin{quote}
\emph{Under mean-field scaling, symmetric fixed-$M$ head diversity is generically a subleading effect unless explicit order-one symmetry breaking is introduced.}
\end{quote}
Therefore, if head diversity remains large in a very wide practical model, at least one of the following must be happening:
\begin{enumerate}[leftmargin=2em]
  \item the training protocol injects order-one asymmetry across heads,
  \item finite-width fluctuations are dynamically amplified,
  \item the network is operating closer to a kernel/lazy regime than to a deterministic mean-field regime.
\end{enumerate}
\end{remark}

\subsection{Implications for TabM and BatchEnsemble}

\begin{remark}[What the theory suggests about efficient ensembles]
\label{rem:practitioner}
The analysis above points to several concrete takeaways for practitioners.
\end{remark}

\paragraph{1. BatchEnsemble can behave more like one shared model than like a true ensemble.}
This is the most direct implication of \cref{cor:collapse,cor:rowvariation,rem:finite-width}.
If the heads are trained under a nearly symmetric protocol, then deterministic feature-learning dynamics pushes them toward a common predictor.
Recent empirical evidence is compatible with this warning:
on image benchmarks, a 2026 study reports that BatchEnsemble tracks a single-model baseline more closely than a deep ensemble, and that members are near-identical in function and parameter space \cite{zamyatin2026}.

\paragraph{2. But collapse is not automatically bad for point prediction.}
Even if the heads collapse, the resulting scalar dynamics is not necessarily the same as that of a plain single MLP.
The effective kernels
\[
  K_{\mathrm{eff}}=\frac{K_{\mathrm d}+(M-1)K_{\mathrm o}}{M}
\]
in the abstract model and
\[
  K_{\mathrm{eff}}^{\mathrm{LevelB}}
  =
  \frac{
    K_d^{\mathrm{neu}}+(M-1)K_o^{\mathrm{neu}}+K^{\mathrm{head}}
  }{M}
\]
in the semifaithful Level-B model still remember the shared/private architecture.
So efficient ensembles may help primarily as
\begin{enumerate}[leftmargin=2em]
  \item structured reparameterizations,
  \item regularizers induced by shared optimization,
  \item capacity-sharing mechanisms that alter the learned feature dynamics.
\end{enumerate}
This viewpoint fits the official TabM guidance that many trained submodels can later be pruned with only a minor performance drop \cite{tabm-repo}.

\paragraph{3. Independent output heads matter, but not in a one-dimensional way.}
The Level-B formulation keeps the output coefficients head-specific while introducing collectivity through shared hidden features and the global modulation variables.
This role is now more nuanced than a simple ``independent heads preserve diversity'' slogan.
\Cref{cor:levelb-shared-collapsed} shows that, in the corresponding shared-output Level-B proxy, the direct last-layer coupling term survives at order one on the symmetric manifold, whereas in the TabM-like Level-B model the corresponding common-mode contribution is averaged down by a factor of $1/M$.
So making the output head independent avoids adding yet another shared route through which heads can couple in the \emph{common predictor}.
However, \cref{prop:levelb-transverse,cor:levelb-mode-allocation} show that independent heads also create extra positive-semidefinite diagonal terms in the local frozen disagreement operator.
So independent heads weaken direct collectivity in the average mode, but they do \emph{not} by themselves guarantee deterministic mean-field diversity.
If the goal is uncertainty or ensemble diversity rather than just point accuracy, independent heads remain a meaningful design choice, but the theory says that one should still expect finite-width effects or explicit symmetry breaking to do much of the practical work.

\paragraph{4. Increasing $k$ without increasing capacity should eventually give diminishing returns.}
The repository notes that, when $k$ grows, one should often increase `d\_block` or `n\_blocks`, precisely because weight sharing makes a larger ensemble harder to accommodate in a fixed backbone \cite{tabm-repo}.
This again matches the theory:
if many heads are forced through the same shared feature bottleneck, then extra heads become increasingly redundant unless
\begin{enumerate}[leftmargin=2em]
  \item the backbone is enlarged,
  \item the sharing pattern is weakened,
  \item or stronger symmetry-breaking is introduced.
\end{enumerate}

\paragraph{5. Domain-specific empirical outcomes are not a contradiction, but a clue.}
Recent evidence is mixed.
The image-based study of \cite{zamyatin2026} argues that BatchEnsemble behaves like a single model for calibration and OOD tasks,
while a separate 2026 study finds that BatchEnsemble can match deep ensembles on tabular and time-series uncertainty estimation tasks \cite{blorstad2026}.
Our reading is that this tension is informative:
whether an efficient ensemble acts like a real ensemble depends on how much non-shared variation survives in a given architecture/data regime.
Tabular models with independent output heads may preserve useful asymmetry better than heavily shared CNN settings.

\paragraph{6. A practical diagnostic checklist.}
If one wants to know whether a BatchEnsemble-like model is functioning as a genuine ensemble or as a shared regularized model, the following ablations are especially informative:
\begin{enumerate}[leftmargin=2em]
  \item compare independent output heads against a shared output layer,
  \item increase $k$ while keeping backbone capacity fixed,
  \item prune trained heads after training and measure the performance drop.
\end{enumerate}
If these interventions barely change either the per-head functions or the ensemble accuracy, then the model is likely operating in the ``effectively one shared predictor'' regime.

\subsection{A focused experimental plan: width and head design}

If the main goal is to test the current theory with the smallest possible experimental surface area,
then the most informative plan is to focus on only two axes:
\begin{enumerate}[leftmargin=2em]
  \item the \emph{output-head design}, namely independent output heads versus a shared output layer,
  \item the \emph{network width}, which controls how close the system is to its infinite-width limit.
\end{enumerate}
Other knobs can be postponed to appendix ablations once these two primary comparisons are understood.

\paragraph{Head axis: independent versus shared output.}
The cleanest architectural experiment is to compare:
\begin{enumerate}[leftmargin=2em]
  \item a TabM-like Level-B model with independent output coefficients $a_{\alpha,j}$,
  \item a matched shared-output proxy in which the last-layer coefficients are shared across heads.
\end{enumerate}
The backbone, optimization hyperparameters, training protocol, and data stream should be matched as closely as possible.
The guiding theoretical prediction is provided by \cref{cor:levelb-shared-collapsed}:
both models may collapse at deterministic leading order,
but the shared-output model carries an order-one direct last-layer collectivity term,
whereas the independent-head contribution is averaged down by a factor of $1/M$.
At the same time, \cref{prop:levelb-transverse,cor:levelb-mode-allocation} show that the independent-head architecture also adds extra diagonal terms to the local frozen disagreement operator.
So the empirical question is not merely whether the heads are identical,
but \emph{which effect dominates at practical width}:
the weaker common-mode collectivity of independent heads,
or the stronger local disagreement contraction of their deterministic transverse operator.

\paragraph{Width axis: finite-width diversity versus deterministic collapse.}
The second core experiment is to sweep width while holding the head design fixed.
This axis is essential because the present theory predicts deterministic collapse only in the infinite-width leading order,
whereas practical diversity may survive as a finite-width effect.
The key qualitative prediction is:
\begin{quote}
\emph{under a symmetric protocol, if head diversity is primarily a finite-width phenomenon,
then per-head disagreement and any ensemble benefit attributable to disagreement should decrease as width grows.}
\end{quote}
The frozen fluctuation template sharpens this qualitative picture:
once the model is wide enough for CLT-scale residual diversity to dominate,
pairwise disagreement should approximately follow an inverse-width law,
equivalently a log-log slope near $-1$.
This is exactly the bridge between the NTK picture and the mean-field picture:
if diversity remains robust even at very large width, then the model is not behaving like the deterministic fixed-$M$ mean-field limit;
if diversity decays with width, then the mean-field collapse story is empirically credible.

\paragraph{Recommended experimental grid.}
A minimal but theory-aligned grid is:
\begin{enumerate}[leftmargin=2em]
  \item two output-head designs: independent versus shared,
  \item at least three or four widths spanning a practically meaningful range,
  \item fixed ensemble size $M$,
  \item matched optimization settings across all runs,
  \item several random seeds, because one wants to separate systematic width effects from seed-to-seed variance.
\end{enumerate}
This already yields a $2\times K$ design that can support the main theoretical claims without introducing too many confounds.

\paragraph{Primary metrics.}
The most informative measurements are:
\begin{enumerate}[leftmargin=2em]
  \item \emph{pairwise head disagreement} on train and test inputs, for example
  \[
    \frac{1}{M(M-1)}
    \sum_{\alpha\neq\beta}
    \mathbb E_x\bigl[(f_\alpha(x)-f_\beta(x))^2\bigr],
  \]
  together with its log-log slope as a function of width,
  \item \emph{pairwise prediction correlation} across heads,
  \item \emph{ensemble gain} over a matched single-head or shared-output baseline,
  \item standard predictive metrics such as accuracy, RMSE, NLL, or calibration error, depending on the task.
\end{enumerate}
The first two metrics measure whether the model behaves like a genuine ensemble.
The third tells us whether any remaining diversity is actually useful.
The fourth checks that diversity is not being increased only by degrading point prediction.

\paragraph{Mode-resolved diagnostics.}
Because the current theory separates common-mode and disagreement-mode effects,
a strong empirical validation should not rely only on final accuracy or final disagreement.
Instead, it should explicitly probe the two modes.
The cleanest plan is to include the following three diagnostics.
\begin{enumerate}[leftmargin=2em]
  \item \emph{Common-mode diagnostic.}
  Start from a deliberately symmetric initialization in which all heads are initialized identically inside each architecture class,
  and train with shared batches.
  Measure the early-time evolution of the mean predictor
  \[
    \bar f_t(x)=\frac1M\sum_{\alpha=1}^M f_{\alpha,t}(x)
  \]
  through quantities such as early training-loss decrease or
  \[
    \mathbb E_x[(\bar f_t(x)-\bar f_0(x))^2].
  \]
  The theoretical prediction is that the shared-output model should exhibit a stronger direct output-level common-mode contribution than the independent-head model,
  because the corresponding output-only term enters with coefficient $1$ rather than $1/M$.

  \item \emph{Transverse diagnostic.}
  Again start from a symmetric initialization, but now inject a tiny zero-sum perturbation across heads,
  for example by adding $\pm \varepsilon$ to matched head-specific output coordinates while keeping the ensemble mean fixed.
  Under shared batches, track the decay of
  \[
    D(t)
    :=
    \frac{1}{M(M-1)}
    \sum_{\alpha\neq\beta}
    \mathbb E_x[(f_{\alpha,t}(x)-f_{\beta,t}(x))^2].
  \]
  A convenient summary is the short-time decay rate
  \[
    \lambda_{\mathrm{disc}}(t_0,t_1)
    :=
    -\frac{1}{t_1-t_0}
    \log\frac{D(t_1)}{D(t_0)}.
  \]
  The frozen-kernel theory predicts that the independent-head model can exhibit \emph{stronger}, not weaker, local disagreement contraction, because of the extra positive-semidefinite terms in \cref{eq:levelb-transverse-ind-kernel}.

  \item \emph{Empirical kernel audit.}
  On a held-out batch near initialization, compute the empirical ensemble tangent kernel blocks
  \[
    \widehat K_{\alpha\beta}(x_i,x_j)
    =
    \langle \nabla_\theta f_\alpha(x_i),\nabla_\theta f_\beta(x_j)\rangle
  \]
  or a tractable approximation restricted to the last layer and global head variables.
  Then compare
  \[
    \widehat K_{\mathrm{avg}}
    :=
    \frac{1}{M^2}\sum_{\alpha,\beta}\widehat K_{\alpha\beta}
    \qquad\text{and}\qquad
    \widehat K_{\mathrm{tr}}
    :=
    \widehat K_{\mathrm d}-\widehat K_{\mathrm o}.
  \]
  This gives the most direct empirical test of the theory:
  shared output should look more collective in the common mode,
  while independent heads should exhibit a larger transverse self-coupling contribution.
\end{enumerate}

\paragraph{A practical three-stage protocol.}
Putting the above together, the cleanest experimental program is:
\begin{enumerate}[leftmargin=2em]
  \item run the \emph{natural-training comparison} across width and head design, using standard random initialization and the shared-batch protocol;
  \item run a \emph{symmetric diagnostic experiment} to isolate common-mode dynamics;
  \item run a \emph{tiny-perturbation experiment} to isolate disagreement-mode contraction.
\end{enumerate}
The first stage tells us what practitioners actually observe.
The second and third stages explain \emph{why} those observations occur in terms of the mode decomposition predicted by the theory.

\paragraph{Expected reading of the results.}
The theory suggests the following interpretation:
\begin{enumerate}[leftmargin=2em]
  \item if shared-output models exhibit lower disagreement and smaller ensemble gain than independent-head models at matched width,
  then beyond-leading-order effects are allowing the weaker common-mode collectivity of independent heads to dominate in practice;
  \item if the disagreement levels are similar, or if independent-head models collapse just as strongly, then the deterministic transverse contraction from \cref{prop:levelb-transverse} is likely the better explanation;
  \item if disagreement decays with width in both models,
  then the deterministic mean-field collapse picture is the right first-order explanation;
  a log-log slope near $-1$ is the most direct support for the frozen CLT template,
  \item if disagreement stays large as width grows, then order-one asymmetry, lazy-regime effects, or finite-width fluctuation amplification must be playing a major role;
  \item if independent heads retain useful diversity at moderate widths while shared-output models do not,
  then the theory has identified a practically meaningful architecture knob, but the mechanism should be interpreted as a beyond-leading-order effect rather than as deterministic local transverse instability.
\end{enumerate}

\paragraph{Why this plan is enough for a first paper.}
This two-axis design is attractive because it is tightly coupled to the current theory.
It directly tests the two strongest practitioner-facing claims already supported by the analysis:
\begin{enumerate}[leftmargin=2em]
  \item independent output heads weaken the most direct route to collectivity,
  \item width controls how strongly the system behaves like its deterministic mean-field limit.
\end{enumerate}
So even before adding auxiliary protocol ablations,
this experimental plan can already validate or falsify the central conceptual story of the paper.

\section{Status of the Current Theory}

At this stage, the results are best organized into three layers:
\begin{enumerate}[leftmargin=2em]
  \item unconditional statements that are fully proved in the characteristic-flow frameworks analyzed here,
  \item exact statements that are either conditional on regular characteristic solutions or concern auxiliary / surrogate models,
  \item statements that remain open and should be presented only as conjectures or next steps.
\end{enumerate}

\subsection{Layer I: unconditional rigorous results}

The following statements are rigorous theorems for the abstract bounded-Lipschitz fixed-$M$ mean-field model from \cref{sec:abstract-fixed-m}:
\begin{enumerate}[leftmargin=2em]
  \item global characteristic well-posedness of the distributional dynamics on every finite time interval (\cref{thm:wellposed}),
  \item stability with respect to the initial law in Wasserstein distance (\cref{prop:stability}),
  \item deterministic particle-to-PDE convergence for empirical measures (\cref{cor:empirical-limit}),
  \item preservation of permutation symmetry across ensemble members (\cref{thm:symmetry}),
  \item exact collapse to a common predictor under symmetric initialization, with the scalar reduced dynamics from \cref{cor:collapse} and effective kernel
  \[
    K_{\mathrm{eff}}=\frac{K_{\mathrm d}+(M-1)K_{\mathrm o}}{M},
  \]
  \item equality of the per-head loss and the loss of the mean prediction on the symmetric manifold (\cref{cor:losses}),
  \item nonlinear stability of the symmetric manifold with respect to the distance from the symmetrized initial law (\cref{thm:nonlinear-stability}),
  \item exact average/deviation decomposition, showing that only centered row-variation of the kernel can drive deterministic disagreement (\cref{prop:deviation,cor:rowvariation}).
\end{enumerate}

For the semifaithful Level-B two-layer model, the following unconditional truncated statements are also rigorous:
\begin{enumerate}[leftmargin=2em]
  \item the truncated coupled PDE-ODE is globally well posed in the characteristic class and preserves permutation symmetry, provided the truncation maps are chosen permutation-equivariantly (\cref{thm:levelb-wellposed}),
  \item the truncated coupled flow is stable with respect to the initial pair $(\rho_0,r_0)$ in the metric $\mathbf d=\W+\|\cdot\|$ (\cref{prop:levelb-stability}),
  \item the truncated finite-width particle/head system converges deterministically to the Level-B limit on every compact time interval (\cref{cor:levelb-empirical-limit}),
  \item around a symmetric Level-B limit, finite-width headwise disagreement vanishes at the deterministic level (\cref{cor:levelb-disagreement}).
\end{enumerate}

\subsection{Layer II: exact but conditional or auxiliary results}

The following statements are mathematically exact, but they either require an a priori regular characteristic solution hypothesis for the untruncated Level-B system or concern auxiliary frozen / surrogate dynamics:
\begin{enumerate}[leftmargin=2em]
  \item conditional on the existence of a regular characteristic solution of the untruncated Level-B system, the output dynamics decomposes into a neuron kernel and an additional head kernel (\cref{prop:levelb-dynamics}),
  \item conditional on the same regularity hypothesis, the collapsed symmetric manifold carries the effective Level-B kernel
  \[
    K_{\mathrm{eff}}^{\mathrm{LevelB}}
    =
    \frac{
      K_d^{\mathrm{neu}}+(M-1)K_o^{\mathrm{neu}}+K^{\mathrm{head}}
    }{M}
  \]
  as stated in \cref{cor:levelb-collapsed},
  \item conditional on the analogous regular characteristic solution hypothesis for the shared-output Level-B proxy, the direct last-layer coupling term survives at order one on the symmetric manifold, in contrast with the $1/M$-suppressed TabM-like independent-head contribution (\cref{cor:levelb-shared-collapsed}),
  \item \Cref{prop:frozen-tangent,cor:frozen-contraction} analyze a frozen-kernel tangent model around a symmetric trajectory.
  They show that average and disagreement modes decouple and that the transverse operator is governed by $K_{\mathrm d}-K_{\mathrm o}$.
  This is a mathematically exact result for the auxiliary linearized system, not for the full feature-learning linearization of the mean-field PDE.
  \item \Cref{prop:levelb-transverse} specializes that frozen transverse operator to the semifaithful Level-B family.
  It shows that independent output heads and shared output differ in opposite ways in the common mode and the disagreement mode:
  independent heads suppress the strongest direct last-layer contribution to the common predictor, but they also add positive-semidefinite diagonal terms to the local disagreement operator.
  \item \Cref{cor:levelb-mode-allocation} packages this into a simple design rule:
  the same output-only feature-product term is allocated differently across modes, appearing with coefficient $1$ in the shared-output common predictor, with coefficient $1/M$ in the independent-head common predictor, with coefficient $0$ in the shared-output disagreement dynamics, and with coefficient $1/M$ in the independent-head disagreement dynamics.
  \item \Cref{prop:frozen-covariance,cor:frozen-fluctuation-decay} give a mathematically exact fluctuation template for the auxiliary frozen tangent system:
  common-mode and disagreement-mode covariances solve separate Lyapunov-type evolution equations, and transverse covariance decays under the same coercivity condition that contracts deterministic disagreement.
  \item \Cref{prop:frozen-pairwise,cor:frozen-pairwise-decay,cor:frozen-width-law} turn that template into directly testable width predictions:
  in the frozen CLT model, pairwise disagreement is exactly proportional to the transverse covariance trace and obeys an exact $m^{-1}$ law.
  \item archived reduced-surrogate and one-feature small-time coupling theorems are now recorded separately in
  \[
    \texttt{tabm\_meanfield\_levelc\_surrogates.tex}.
  \]
  They remain useful as mechanism calculations, but they are no longer part of the main Level-B note.
\end{enumerate}

\subsection{Layer III: open statements}

The following statements have \emph{not} been proved in this note and should currently be treated as open problems:
\begin{enumerate}[leftmargin=2em]
  \item a finite-horizon untruncated Level-B well-posedness / stability / deterministic particle-limit theorem in a moment-controlled $\PP_2(\Theta)\times \mathcal R$ framework, assuming essentially bounded initial head coordinates and finite second moments in $(w,b)$,
  \item an extension of the unconditional characteristic-flow theory beyond uniformly bounded loss derivatives, for example to squared loss via moment-controlled estimates,
  \item a full propagation-of-chaos theorem for an architecture-faithful layerwise BatchEnsemble/TabM surrogate,
  \item a theorem that the Level-B model or an architecture-faithful layerwise surrogate necessarily creates off-diagonal coupling at positive times,
  \item a fluctuation or central-limit theorem explaining finite-width diversity around the deterministic symmetric solution.
\end{enumerate}

Among these, the first item now looks more like a tractable technical upgrade than a conceptual roadblock.
The present drift structure indicates that the real obstruction is the mixed observable $a_\alpha w$ in the head ODE:
it breaks the bounded-Lipschitz $\W+\|\cdot\|$ argument used in
\cref{thm:levelb-wellposed,prop:levelb-stability},
but it should be replaceable by finite-horizon bounds on $(a_\alpha)_\alpha$,
propagation of first and second moments in $(w,b)$,
and a local contraction on bounded-moment sets in the metric $W_2+\|\cdot\|$.
This heuristic is consistent with the role of coefficient-growth control in two-layer mean-field neural-network theory \cite{mei2019}
and with locally Lipschitz truncation/continuation arguments for McKean--Vlasov equations \cite{erny2022};
see also \cite{chizat2018} for the ambient measure-gradient-flow viewpoint.

\subsection{Takeaway}

The cleanest safe summary is therefore the following:
\begin{quote}
\emph{At deterministic fixed-$M$ mean-field level, the bounded-Lipschitz characteristic-flow theory already forces symmetric BatchEnsemble/TabM-like dynamics to collapse to a common predictor and remain quantitatively stable under small asymmetry.
In the semifaithful Level-B model, the truncated particle/head system already converges deterministically to a coupled PDE-ODE limit, and for regular characteristic solutions of the untruncated system the collapsed dynamics carries an additional head-kernel contribution coming from the global modulation variables.
Within the same architecture family, this stands in genuine contrast to NTK/GP analyses, where order-one initialization randomness can keep an efficient ensemble in an independent regime.
The shared-output versus independent-head comparison is also genuinely two-sided:
independent heads weaken direct common-mode collectivity, but in the local frozen transverse operator they contribute additional positive-semidefinite self-coupling terms rather than creating deterministic mean-field diversity on their own.
The frozen tangent system already supplies a clean covariance template for any future CLT-scale fluctuation theory, but a full finite-width fluctuation theorem explaining practical diversity remains open.}
\end{quote}

Those missing pieces are natural next steps.
But the current note already isolates a mathematically clean backbone that can support a stronger later paper.

\section{Recommended Next Theorem}

The most natural next rigorous step is no longer a vague ``remove the truncation somehow,'' but the following finite-horizon upgrade of Level B:

\begin{quote}
\emph{Assume bounded inputs, $\sigma\in C_b^2(\R)$, and a globally Lipschitz uniformly bounded loss derivative $\partial_1\ell(\cdot,y)$.
Assume also that $\rho_0\in \PP_2(\Theta)$, that the output-head coordinates $(a_\alpha)_\alpha$ are essentially bounded under $\rho_0$, and that $r_0\in\mathcal R$.
Then for every $T>0$, the untruncated Level-B system admits a unique solution
\[
  (\rho_t,r_t)_{t\in[0,T]}
  \in
  C\bigl([0,T],\PP_2(\Theta)\times\mathcal R\bigr).
\]
Moreover, the coordinates $(a_\alpha(t))_\alpha$ remain uniformly bounded on $[0,T]$,
the second moments of $(w,b)$ remain finite on $[0,T]$,
permutation symmetry is preserved,
and both stability and deterministic particle-to-limit convergence hold on $[0,T]$ in a moment-aware metric such as $W_2+\|\cdot\|$.}
\end{quote}

This is the shortest realistic path to turning \cref{prop:levelb-dynamics,cor:levelb-collapsed}
from conditional regular-characteristic statements into consequences of a fully rigorous untruncated Level-B theory.
After that upgrade, the next genuinely new difficulty is no longer truncation removal itself,
but passing from Level B to an architecture-faithful layerwise surrogate without losing control of the shared head variables.

\begin{thebibliography}{9}

\bibitem{tabm}
Y.~Gorishniy, A.~Kotelnikov, and A.~Babenko.
\newblock \emph{TabM: Advancing Tabular Deep Learning with Parameter-Efficient Ensembling}.
\newblock arXiv:2410.24210, 2024.

\bibitem{tabm-repo}
Y.~Gorishniy, A.~Kotelnikov, and A.~Babenko.
\newblock \emph{Official TabM repository}.
\newblock \url{https://github.com/yandex-research/tabm}, accessed 2026-03-23.

\bibitem{batchensemble}
Y.~Wen, D.~Tran, and J.~Ba.
\newblock \emph{BatchEnsemble: An Alternative Approach to Efficient Ensemble and Lifelong Learning}.
\newblock ICLR, 2020.

\bibitem{embedded-ensembles}
M.~Velikanov, R.~Kail, I.~Anokhin, R.~Vashurin, M.~Panov, A.~Zaytsev, and D.~Yarotsky.
\newblock \emph{Embedded Ensembles: Infinite Width Limit and Operating Regimes}.
\newblock AISTATS, 2022.

\bibitem{mei2018}
S.~Mei, A.~Montanari, and P.-M.~Nguyen.
\newblock \emph{A Mean Field View of the Landscape of Two-Layer Neural Networks}.
\newblock Proceedings of the National Academy of Sciences, 115(33):E7665--E7671, 2018.

\bibitem{mei2019}
S.~Mei, T.~Misiakiewicz, and A.~Montanari.
\newblock \emph{Mean-Field Theory of Two-Layers Neural Networks: Dimension-Free Bounds and Kernel Limit}.
\newblock COLT, 2019.

\bibitem{chizat2018}
L.~Chizat and F.~Bach.
\newblock \emph{On the Global Convergence of Gradient Descent for Over-parameterized Models using Optimal Transport}.
\newblock NeurIPS, 2018.

\bibitem{erny2022}
X.~Erny.
\newblock \emph{Well-posedness and Propagation of Chaos for McKean--Vlasov Equations with Jumps and Locally Lipschitz Coefficients}.
\newblock Stochastic Processes and their Applications, 150:192--214, 2022.

\bibitem{zamyatin2026}
A.~Zamyatin, P.~Indri, S.~Malhotra, and T.~G\"artner.
\newblock \emph{Is BatchEnsemble a Single Model? On Calibration and Diversity of Efficient Ensembles}.
\newblock arXiv:2601.16936, 2026.

\bibitem{blorstad2026}
M.~Bl{\o}rstad, H.~J.~Mostein, N.~Blaser, and P.~Parviainen.
\newblock \emph{Evaluating Prediction Uncertainty Estimates from BatchEnsemble}.
\newblock arXiv:2601.21581, 2026.

\end{thebibliography}

\end{document}
